\documentclass[11pt]{article}

\usepackage{fontspec}
\defaultfontfeatures{Renderer=Harfbuzz,Ligatures=TeX}
\setmainfont{Noto Sans}
\setsansfont{Noto Sans}
\setmonofont{Noto Sans Mono}
\usepackage[margin=1in]{geometry}
\usepackage{amsmath,amssymb}
\usepackage{booktabs}
\usepackage{array}
\usepackage{longtable}
\usepackage{tabularx}
\usepackage{graphicx}
\usepackage{caption}
\usepackage{enumitem}
\usepackage{ragged2e}
\usepackage{hyperref}
\usepackage{url}
\usepackage{polyglossia}
\setdefaultlanguage{serbian}
\newfontfamily\cyrillicfont[Script=Cyrillic]{Noto Sans}
\newcolumntype{P}[1]{>{\raggedright\arraybackslash}p{#1}}
\captionsetup{font=small,labelfont=bf}
\setlength{\emergencystretch}{6em}
\hypersetup{
  unicode=true,
  colorlinks=true,
  linkcolor=blue,
  citecolor=blue,
  urlcolor=blue,
  pdflang={sr},
  pdftitle={Cosmo-Local Credit (CLC) White Paper v0.8},
  pdfauthor={William O. Ruddick and Mohamed Sohail}
}

\title{Cosmo-Local Credit (CLC)\\Мрежа за усмеравање кредита, усклађивање обавеза и финансирање здраве космо-локалне економије}
\author{William O. Ruddick \\ Mohamed Sohail \\ Grassroots Economics Foundation \\ \texttt{info@grassecon.org}}
\date{White Paper v0.8 translation: 10 October 2026}

\begin{document}
\maketitle
\tableofcontents
\newpage
\sloppy
\RaggedRight

\textbf{О овом преводу.} Ово је превод White Paper v0.8. Енглески изворни текст објављен је 30. септембра 2026. Овај превод је објављен 10. октобра 2026. Енглеска верзија је изворни текст. \href{https://docs.cosmolocal.credit/white-paper}{Прочитајте енглески извор}.

\textbf{Верзија:} 0.8

\textbf{Датум:} 30. септембар 2026.

\textbf{Аутори:} William O. Ruddick и Mohamed Sohail

\textbf{Контакт:} \href{mailto:info@grassecon.org}{info@grassecon.org}

\textbf{Преузимање:} \href{https://docs.cosmolocal.credit/white-paper/Cosmo-Local-Credit-CLC-White-Paper-v8-sr.pdf}{CLC White Paper v0.8 PDF на српском}

\textbf{Издање:} Коначно јавно издање са верзијом. Овај документ није инвестициони савет и не представља понуду за продају нити позив за куповину хартије од вредности или финансијског инструмента. Касније верзије могу изменити овде описане моделе.

\section*{Како читати овај документ}
\addcontentsline{toc}{section}{Како читати овај документ}

Документ обједињује четири врсте садржаја. Статус важи у целом тексту чак и када га поглавље не понавља.

\begin{longtable}{P{0.30\linewidth} P{0.62\linewidth}}
\toprule
Статус & Значење \\
\midrule
\endhead
\textbf{Актуелно} & Понашање потврђено у CLC App или јавним уговорима Protocol v1.1.0. \href{https://docs.cosmolocal.credit/sr/protocol/overview}{Референца Протокола} је чињенични извор за те уговоре. \\
\textbf{Зависно од примене} & Могућност која постоји само тамо где је омогуће конкретна Тржница, пружалац услуге, јурисдикција, орган управљања или објављена политика. \\
\textbf{Историјско} & Доказ или могућност раније услуге Sarafu Network. Не представља актуелну употребу CLC-а нити доказ да су корисник, имовина, запис или обавеза пренети. \\
\textbf{Предложено} & Модел, економски оквир, механизам управљања или ставка плана развоја која није представљена као примењена. \\
\bottomrule
\end{longtable}

Актуелни Protocol v1.1.0 обухвата непосредно извршење у \texttt{SwapPool}, необавезне модуле избора и вредновања, ограничења стања токена у Тржници, накнаде Тржнице, додатну протоколску накнаду и \texttt{SwapRouter} који само даје понуде. Вишекорачно извршење, HTLC или ескроу усмеравање, групно нетирање, заједничко осигурање, предложени CLC токен управљања, предложена CLC Мрежна тржница, механизми кредита за накнаде и програми мрежне ликвидности предложени су или зависе од примене.

\textbf{Намењено:} градитељима и Управницима Тржница, инжењерима и ревизорима протокола, финансијерима ликвидности и мандата, пројектантима управљања и правним стручњацима или стручњацима за усклађеност.

\section*{Сажетак}
\addcontentsline{toc}{section}{Сажетак}

Cosmo-Local Credit је породица производа заснована на \textbf{Протоколу удруживања обавеза (CPP)}. CPP описује како независно управљане Тржнице обавеза могу бирати обавезе које се могу искористити, објавити методе курса и ограничења, држати залиху и омогућити одговорну размену.

\textbf{Тржница обавеза} је уређени систем којим се управља. Актуелни \textbf{\texttt{SwapPool}} у Protocol v1.1.0 је једна компонента уговора: трезор токена и механизам непосредне размене. Примена може повезати тај уговор са регистрима, \texttt{Quoter} уговорима, ограничењима стања токена у Тржници, политикама накнада и другим услугама.

Космо-локално значи делити отворене стандарде, софтвер и знање глобално, а издавање, испуњење, управљање и стварну одговорност задржати локално. Издавалац остаје одговоран за свој ваучер. Управник Тржнице остаје одговоран за правила Тржнице и сваку гаранцију коју изричито прихвати. Ни CLC App ни GEF не гарантују аутоматски ваучер, Тржницу, вредност, путању, положај ликвидности нити исход.

\section*{Актуелна примена}
\addcontentsline{toc}{section}{Актуелна примена}

GEF управља CLC App на \href{https://cosmolocal.credit}{cosmolocal.credit}. Она нуди приступ налогу и новчанику, јавно Тржиште и подржане интерфејсе за токене, ваучере, Понуде, Тржнице обавеза, преносе и непосредне размене преко Тржница. Доступност зависи од интерфејса, уговора, залихе, подешених ограничења и накнада, стања мреже, услова приступа, пружалаца услуга, јурисдикције и објављених услова издаваоца или Тржнице.

Само управљање интерфејсом не чини GEF издаваоцем, Управником Тржнице, чуваром имовине, зајмодавцем, зајмопримцем, посредником, гарантом, страном која искоришћава ваучер, саветником, осигуравачем нити страном у обавезама између учесника. Употребу апликације уређују \href{https://docs.cosmolocal.credit/sr/governance/terms}{Услови коришћења}.

Protocol v1.1.0 раздваја одговорне улоге од техничке контроле. Управник Тржнице, власник \texttt{SwapPool}, администратор проксија, контролор зависности, модератор каталога и прималац накнаде могу бити различите стране. За заједничку терминологију погледајте \href{https://docs.cosmolocal.credit/sr/introduction/concepts}{Појмове и речник}.

\section*{Историјско порекло}
\addcontentsline{toc}{section}{Историјско порекло}

Историјска Sarafu Network претходила је CLC App и пружила доказе о валутама заједнице и удруживању обавеза. Прелаз услуге на Cosmo-Local Credit није пренео сваки историјски налог, новчаник, токен, ваучер, Тржницу, стање, извештај, учесника нити обавезу.

Историјски показатељи у овом документу користе датиран и унутрашње доследан скуп података Sarafu. То нису актуелни показатељи употребе CLC-а. Извор и период мерења наведени су у \href{https://docs.cosmolocal.credit/sr/white-paper/executive-summary}{Сажетку}.

\section*{Модел дизајна}
\addcontentsline{toc}{section}{Модел дизајна}

CPP користи четири широке функције:

\begin{itemize}[leftmargin=*]
\item \textbf{Избор:} одређивање ваучера или друге имовине која се прихвата.
\item \textbf{Вредновање:} објављивање метода за стварање курса или понуда Тржнице.
\item \textbf{Ограничавање:} примена актуелних ограничења стања токена у Тржници или других засебно примењених контрола ризика.
\item \textbf{Размена:} држање залихе и извршавање размена преко Тржнице уз објављене накнаде и границе.
\end{itemize}
Примена може додати усмеравање, праћење, подршку ликвидности, гаранције, резерве, осигурање или услуге управљања само према засебно објављеним политикама. Приказ, избор, ограничење, резерва или блокчејн запис сами по себи не утврђују капацитет издаваоца, испуњење, препоруку, гаранцију нити осигурање.

\section*{Вредности и начела дизајна}
\addcontentsline{toc}{section}{Вредности и начела дизајна}

\begin{itemize}[leftmargin=*]
\item \textbf{Брига о људима:} подршка појединачном и заједничком благостању.
\item \textbf{Брига о животној средини:} смањење еколошке штете и подршка регенеративној пракси.
\item \textbf{Правичност:} подстицање безбедног и праведног приступа ресурсима, знању и бризи.
\item \textbf{Узајамност:} дељење ризика, трошкова, одговорности и вишка према објављеним правилима.
\item \textbf{Недоминација:} отпор концентрисаној контроли над људима, подацима, финансијама, знањем и материјалним ресурсима.
\item \textbf{Отпорност:} помоћ заједницама да се припреме и прилагоде економским, политичким и климатским поремећајима.
\item \textbf{Локална одговорност:} задржавање стварног испуњења и правне одговорности код наведених страна.
\item \textbf{Могућност издвајања:} омогућавање усаглашеним заједницама и применама да напусте заједничке регистре или услуге без брисања важећих стања или обавеза.
\end{itemize}
Дигитални слој је заједничко памћење и инфраструктура координације. Није замена за односе, испуњење издаваоца, локално управљање нити правну одговорност.

\section*{Актуелни и предложени токови}
\addcontentsline{toc}{section}{Актуелни и предложени токови}

\subsection*{Актуелни непосредни ток Тржнице}
\addcontentsline{toc}{subsection}{Актуелни непосредни ток Тржнице}

\begin{enumerate}[leftmargin=*]
\item Ималац проверава ваучер, издаваоца, услове и повезане Понуде.
\item Тржница прихвата подржану имовину и објављује подешавања и одговорне органе.
\item Ималац добија понуду и одобрава непосредну размену преко Тржнице.
\item \texttt{SwapPool} врши поравнање размене на блокчејну и обрачунава накнаду Тржнице и протокола.
\item Ако ималац касније изабере \textbf{„Искористи``} у апликацији, апликација припрема пренос власнику токена као подношење ради искоришћења.
\item Издавалац засебно извршава обећано испуњење и бележи раздужење како се испуњене јединице не би поново користиле.
\end{enumerate}
Узимање залихе из Тржнице је размена. Представља искоришћење само ако је Тржница засебно прихватила улогу издаваоца и испуњава примењиве услове ваучера.

\subsection*{Шири предложени модел}
\addcontentsline{toc}{subsection}{Шири предложени модел}

Предложени модел истражује усмеравање извршења, нетирање, заједничке услуге, програме мрежне ликвидности, полисе осигурања и средства управљања. Сваки од ових елемената захтевао би примену, одговорне стране, усвојена правила, финансирање, контроле и јавне објаве пре него што почне да важи.

Предложени модел накнада може користити удео мреже из накнада Тржница и засебне накнаде за усмеравање или услуге. Тај предлог се разликује од Protocol v1.1.0, у којем је необавезна протоколска накнада додатна накнади Тржнице.

Подржаваоци или пружаоци ликвидности не добијају аутоматски уделе у Тржници, отплату, право на повлачење, награду или право управљања из актуелног \texttt{SwapPool}. Сваки актуелни или будући програм мора засебно објавити врсту преноса, одговорну страну, контролу, права на опоравак или повлачење, губитке, награде, накнаде и права управљања.

\section*{Водич за читаоце}
\addcontentsline{toc}{section}{Водич за читаоце}

\begin{itemize}[leftmargin=*]
\item Актуелна терминологија и животни циклус радњи: \href{https://docs.cosmolocal.credit/sr/introduction/concepts}{Појмови и речник}
\item Проверени примењени уговори: \href{https://docs.cosmolocal.credit/sr/protocol/overview}{Protocol v1.1.0}
\item Историјски докази: \href{https://docs.cosmolocal.credit/sr/white-paper/executive-summary}{Сажетак}
\item Федерација и интероперабилност: поглавља 5, 8 и 11
\item Предложене гаранције и границе осигурања: поглавља 10 и 11
\item Предложени модели ликвидности и накнада: поглавља 7, 9 и 12
\item Предложени оквир мерења: поглавље 14 и додаци A--C
\item Предложени план развоја: поглавље 15
\item Правна питања и усклађеност: поглавље 17
\end{itemize}
\href{https://docs.cosmolocal.credit/sr/white-paper/archive}{Претходне верзије Беле књиге} задржане су само као јасно означене замењене историјске публикације.

\section*{Сажетак}
\addcontentsline{toc}{section}{Сажетак}

У овом случају, \textbf{Протокол удруживања обавеза (CPP)} описује како Тржнице обавеза, који је независно регулисан, може да контролише обавезе које се могу откупити, објављује методе и границе размене, води инвентар и омогућава одговорну размену. Емитенти остају одговорни за своје обавезе по ваучеру. Управници Тржница остаје одговорна за правила групе и сва гаранције које они изричито преузимају. Ни CLC App ни GEF нису универзални гаранти.

Protocol v1.1.0 обезбеђује тренутне темељне блокове директне размене: \texttt{GiftableToken}, \texttt{SwapPool}, изборне регистре и модуле за процјену, ограничења баланса токена пула, таксе пула, додатна такса протокола и само цитат \texttt{SwapRouter}. Садашњи договори не записују реализацију у стварном свету, не стварају аутоматске акције пружаоца ликвидности, не извршавају маршруте са више хопа, нити пружају универзалне резерве, гаранције или осигурање.

Историјски \textbf{Sarafu Network} користили су концепте заједничког ангажовања у заједничким валутама, групама за штедњу, системима заједничке помоћи и обавезама производње. Његова служба је касније прешла на CLC App. Ово није било представљање да је сваки историјски рачун, новчаник, токен, ваучер, пул, баланс, извештај, учесник или обавеза мигрирао.

\subsection*{Историјска снимка Sarafu Network --- Активност Celo од 5. јула 2023. до 20. јула 2025.}
\addcontentsline{toc}{subsection}{Историјска снимка Sarafu Network --- Активност Celo од 5. јула 2023. до 20. јула 2025.}

\textbf{Извор:} \href{https://dune.com/grassrootseconomics/sarafu-network}{Дашборд Дјуна Аналитика}

\begin{itemize}[leftmargin=*]
\item 26.367 корисника
\item 285.197 размене међусобног броја
\item 188 јединствена активна Тржнице обавеза
\item 745 јединствене активне ваучере
\item \$320,692 Област суапа пула
\item 899 објављени извештаји (\href{https://sarafu.network/reports}{архив историјских извештаја})
\end{itemize}
Ове бројке су историјски докази, а не тренутне CLC бројке употребе.

Широки предложен дизајн CLC истражује како су компатибилни мрежи могли:

\begin{enumerate}[leftmargin=*]
\item открити и цитирати путеве преко независно управљаних Тржница;
\item Координира одговорне мрежне услуге без преоптерећења локалне одговорности;
\item подршка одвојено финансираним мандатима ликвидности, гаранцијама, резервима или осигурањима;
\item мерења суапа група одвојено од презентације ваучера, испуњавања и ослобођења; и
\item коришћење предложених мрежних трошкова и трошкова за услуге за финансирање усвојених заједничких услуга.
\end{enumerate}
Предлог обухвата \textbf{предложени CLC токен управљања} и \textbf{предложену CLC Мрежну тржницу}. Ни једно ни друго није део актуелне апликације или Protocol v1.1.0.

Према предложеном моделу, учесници у ликвидности и управљању могли су да делују само путем одвојеног усвојења програма са објављеним условима прихватљивости, ризицима, контролама, правама преноса, условима повлачења и правилама плате. Обични депозити у пулу не стварају аутоматске акције, повраћај, повраћај, награду или право на управљање.

Спољни циљ је:

Повећати одговорну размену и испуњавање обавеза у стварном свету, сачувајући бригу, праведност, локалну одговорност и отпорност.

Протичавање заробљавања и веродостојни излазак остају основни пројектни циљеви. Предложену контролу укључују управљање одложеном временом, вишеструке прагове одобрења, транспарентне делегације, правила конфликта, преглед инцидента и документовани процес раздвајања и миграције. Ове контроле би смањиле одређене ризике управљања, али не би могли елиминисати губитак или гарантовати резултате.


\section*{1. Протокол о уједињењу обавеза (CPP): основна примитивна}
\addcontentsline{toc}{section}{1. Протокол о уједињењу обавеза (CPP): основна примитивна}

\textbf{Психијски модел:} Тржница обавеза је уређење које се регулише за прихватање купонова или других средстава, објављивање правила размене, држање инвентаризације и омогућивање суапа. Власници касније испоручују ваучере својим емитентима како би се испунили у стварном свету. Измена пула и испуњавање емитента су одвојени животни циклуси.

CPP координише вредност кроз јасно описане обавезе. Модел се разматра у \href{https://willruddick.substack.com/p/grassroots-economics-the-book-is}{Граничне економије: размишљање и пракса}.

\subsection*{1.1 Шта је посвећеност?}
\addcontentsline{toc}{subsection}{1.1 Шта је посвећеност?}

Одговор је обећање идентификоване стране да ће у будућности доставити, на пример, храну, транспорт, рад, складиштење или другу законску робу, услугу, корист или износ. А \textbf{купон} је знак или запис представљен као таква обавеза у објављеним условима.

Токенски договор снима дигиталну механику. Услови купона идентификују емитента, понуду, капацитет, место, време, ограничења, презентацију, испуњавање, жалбе и процес ослобођења.

\subsection*{1.2 Шта је Тржница обавеза?}
\addcontentsline{toc}{subsection}{1.2 Шта је Тржница обавеза?}

Тржница обавеза је уређење које се регулише. Она може да управља појединца, кооперација, заједничка група, јавна агенција, федерација, мултисиг, оператор услуга или друга одговорна структура.

Одлучне улоге и овлашћења укључују:

\begin{itemize}[leftmargin=*]
\item \textbf{Управник Тржнице:} Публикује и управља правилима групе и свим изричито преузиманим гаранцијама;
\item \textbf{Собственик Тржница:} има тренутне овлашћења власника \texttt{SwapPool};
\item \textbf{прокси администратор:} могу да унапреде прокси имплементацију;
\item \textbf{контролери зависности:} управљају конфигурисаним регистровима, котирама, ограничивачима или компонентима плате;
\item \textbf{тражиоц или оператор пута:} могу да открију цитати или, у будућој имплементацији, извршавају одвојен овлашћен пут; и
\item \textbf{гарант:} преузима одређену обавезу само путем објављених, финансираних услова.
\end{itemize}
групе CPP Функције пула су подељене на четири концепта:

\begin{itemize}[leftmargin=*]
\item \textbf{Курација:} прихватају подржане токове или купоне.
\item \textbf{Оценка:} објављују методу коришћену за девизни курс или цитат.
\item \textbf{Ограничење:} примењују тренутне границе баланса токена пула или друге посебно имплементисане контроле.
\item \textbf{Размена:} држе инвентар, извршавају суапе, рачунају таксе и издају записи трансакција.
\end{itemize}
Protocol v1.1.0 имплементира ове функције кроз \texttt{SwapPool} и опционалне зависности. Његов тренутни \texttt{Limiter} ограничава токенски баланс у пулу; не предвиђа ограничења за ролинг, по рачуну или по мрежи. Његов тренутни \texttt{SwapRouter} рачуна котирања за више пула; не извршава суапе.

Широки предложен дизајн CPP може додати ограничења за ролинг, контроле рачуна, рутере за извршење, HTLC или траке за гаранцију и мрежу за партије. То су предложене компоненте, а не описи стручног ограничивача или рутера.

\subsection*{1.3 Логика садашњег директног обмена}
\addcontentsline{toc}{subsection}{1.3 Логика садашњег директног обмена}

Тренутна директна суап пула:

\begin{enumerate}[leftmargin=*]
\item Проверава изборни регистар улазних и излазних токова;
\item мере примљену информацију;
\item добија цитат из конфигурисаног цитирача или примењује парату сировине јединице;
\item Проверка резултирајућег баланса токена пула против опционалног ограничивача;
\item израчује таксу групе и било какву додатну таксу протокола;
\item проверава доступну производну листу;
\item преноси протоколску таксу и излаз и рачуна за таксу групе; и
\item емитује догађаје суапа.
\end{enumerate}
Цитат је параметр трансакције, а не доказ капацитета емитента, вредности откупа, праведности, конвертибилности новца или гаранције.

\subsection*{1.4 Шире употребе}
\addcontentsline{toc}{subsection}{1.4 Шире употребе}

Тржнице обавеза може подржати заједничку размену, производњу, међусобну помоћ, јавне програме и друге структуре одговорне. Одвојено документовани кредитни производ може користити ваучер као обезпеку или инструмент за исплату, али би то захтевало додатне и специфичне за трансакцију услове. Обична пошаља, депозит, суап, презентација изкупа, испуњење или испуњење не представљају аутоматски кредит или повраћај.

CPP је намењен за одговорну размену и аудитне трансакције, а не за спекулативни хурн. Записи на ланцу још увек не доказују реализацију у стварном свету или социјални утицај.


\section*{2. Стварање рачуноводства: од актива у доверење}
\addcontentsline{toc}{section}{2. Стварање рачуноводства: од актива у доверење}

Традиционалне финансије полазе од \textbf{Имовина - Обавезе = Капитал}. CPP то преобликује за економију обавеза:

\begin{itemize}[leftmargin=*]
\item \textbf{Способност прихватања:} колико је пула или мрежа још увек спремна да прихвати обавезе емитента.
\item \textbf{Остале обавезе:} издали обећања које остају неуповршене и држе други.
\item \textbf{Капацитет испуњавања:} доказана способност емитента да испуни те обећања.
\end{itemize}
\subsection*{Илустративни однос обавеза и капацитета:}
\addcontentsline{toc}{subsection}{Илустративни однос обавеза и капацитета:}

Способност прихватања - непоплаћене обавезе = преостала способност прихватања

Овај концептуелни идентитет може да информише границе ризика полу: неограничено издање не подразумева неограничено прихватање. То није идентитет баланса или изјава да је сваки ваучер законски дуг или кредит.

\textbf{Пример:} Транспортер издаје 100100 путовања'' у ваучерима. Тржница прихвата максимум 40 путовања одједном. Уколико је 15 прихваћеног путовања остало неизвесно, Тржница има способност да прихвате до 25 додатних путовања у оквиру ове политике. Рачун сам по себи не ствара кредит или испуњење гаранције.



\section*{3. Активност испуњавања, ослобођења и размене}
\addcontentsline{toc}{section}{3. Активност испуњавања, ослобођења и размене}

Промене супона које адресују посебне имовине. Површавање ваучера се дешава када емитент испуни објављену обавезу. Увођење записује испуњене јединице тако да се не могу поново представити. Ови догађаји се не могу раздвојити у једну употребу сељења.

Предложени оквир мерења користи одвојене податке за:

\begin{itemize}[leftmargin=*]
\item завршен обем суапа резервоара;
\item важећи подаци о откупу;
\item потврђена испуњење емитента;
\item ослобођење;
\item неплаћени прихватљиви обавези; и
\item Измерени инвентар Пола.
\end{itemize}
Предложена брзина испуњавања обавезе може упоредити испуњене вриједности током периода са просечној допустивој вредности обавезе, али само када обе користе исте методе оповестене вредности. Доставка токова није довољан знаменац: може укључивати инвентар емитента, истекале или недоступне јединице, тестове и балансе који не представљају неиздатну обавезу треће стране.

Одвојен однос суапа-активности пула може упоредити завршну вредност суапа пула са мерењем инвентарства пула. Он описује девизне активности, а не испуњавање емитента.

Ликвидност може побољшати доступност инвентара и олакшати размене. То не гарантује да ће емитенти остварити резултате, да ће доприносиоци поврати средства или да котирање полу представља изкупу или новчану вредност. Свака права пружаоца ликвидности захтевају одвојене објављене услове.

Видите ли \href{https://docs.cosmolocal.credit/sr/white-paper/appendix-a-math-box}{Додатко А} за дефиниције и \href{https://docs.cosmolocal.credit/sr/white-paper/appendix-c-kpi-definitions}{Додатко С} за предложену спецификацију мерења.


\section*{4. Поново употребљиво обезбеђење налик форварду}
\addcontentsline{toc}{section}{4. Поново употребљиво обезбеђење налик форварду}

Различно документовано кредитно средство може прихватити фунгибилне, преносиве ваучере као обезпеку или као инструмент за исплату. Ако је било:

\begin{itemize}[leftmargin=*]
\item Пријем резерви може учинити обезпеку више откривеном;
\item Додатне опције законског представљања и испуњавања могу подржати повраћај;
\item Лимити, инвентар, процене, ликвидности, перформансе емитента и ризици од неизплате би остали; и
\item ниједан маршрут, испуњење, повратак, вредност или повратак не би били гарантовани.
\end{itemize}
\subsection*{4.1 Кредитна круга произвођача}
\addcontentsline{toc}{subsection}{4.1 Кредитна круга произвођача}

Будућа CLC- компатибилна услуга могла би да подржи кредитирање произвођача само када су стране успоставиле посебан механизам и изрично представиле трансакцију као исплату аванса. У условима трансакције би се идентификовао кредитор и дужничар и навео:

\begin{itemize}[leftmargin=*]
\item суми аванса и изплате;
\item дате истека, казни, таксе и дозвољено коришћење обезпеке;
\item доделу и било какве промене кредитора или носиоца захтева;
\item како се испуњавање ваучера процењује и доказује за рачуноводство повратака;
\item делимични повраћај, преостале суме, неизплата и ослобођење; и
\item доступна средства за опоравак у складу са применитим законом.
\end{itemize}
Илустративни поток је:

\begin{enumerate}[leftmargin=*]
\item Произвођач или пружаоц услуга издаје ваучер који представља обавезу на будућу производњу, као што су десет такси, 50 килограма кукурузе или десет радних сати.
\item Управник Тржнице признаје тај ваучер и објављује методу обмене стопе, границе, накнаде, правила за залихи и било коју изрично претпостављену заштиту.
\item Кредитор или програм ликвидности одвојено обезбеђује аванс у складу са условима писмених трансакција и прихвата купонове произвођача као обезпеку или договорени инструмент повратака.
\item Власници могу да пренесу ваучере, размене их преко компатибилних пулова или да их представију емитенту.
\item Потврђено испуњење емитента задовољава обавезу о ваучеру. Она смањује посебан кредитни баланс само у мери, у вредности и на доказу наведеном у кредитним условима.
\item У документима се разликује испуњење ваучера од рачуноводства кредита и идентификује било који делимични повраћај, остатак износа, доделу, неизплату или ослобођење.
\end{enumerate}
Ова структура би могла омогућити уговорени повраћај у природи, сачувајући одвојене записи за ваучер и кредитне обавезе. То не произлази из обичног преноса, депозита, суапа резерви, презентације откупа, испуњавања или ослобођења.

\textbf{На пример:} Произвођач добија аванс од 1.000 долара под условима у којима се потврђено испуњење одређених купонова за кукурузу кредитује према износу повраћања користећи наведену методу процене. Кредит се користи само након што добије доказе потребне у овим условима. Ако услови кредита не чине ту повезаност, добијање, пренос, размена, представљање или испуњавање ваучера не утиче на аванс.


\section*{5. Од изолованих Тржница до повезане мреже}
\addcontentsline{toc}{section}{5. Од изолованих Тржница до повезане мреже}

Тренутни Protocol v1.1.0 подржава директно извршење преко једног \texttt{SwapPool} и пружа цитат само \texttt{SwapRouter}. Она не врши више-хоп маршруте, ХТЛЦ-а, маршруте за подношење гаранције, сетну мрежу или клиринг преко мреже.

У овом поглављу се предлаже како би независно управљане базе могли да се координирају без да се одустају од сопствених правила прихватања, процене, границе, таксе, инвентаризације, овлашћења и управљања.

\subsection*{5.1 Одвојене мере размене и испуњавања}
\addcontentsline{toc}{subsection}{5.1 Одвојене мере размене и испуњавања}

Федерација би могла побољшати приступ инвентаризацији, али не би спојила ваучер и разменила животни циклус. Свака имплементација би мерела ове догађаје одвојено:

\begin{enumerate}[leftmargin=*]
\item цитирана је рута;
\item један или више суапа покупности извршавају и уређују на ланцу;
\item држављан представи емитенту ваучерске јединице;
\item емитент испуњава обавезу; и
\item испуњене јединице се испуњавају.
\end{enumerate}
Више цитираних или извршених путева не доказују више испуњавања. У извештајима би се навела кохорта, период, средства, метода процене и временски печат, искључења, исправке и докази изван ланца који се захтевају у Додатку Ц.

\textbf{Илустративни пут:} Школа има купоне за кукурузу, али је потребна транспортна купона. Служба маршрута идентификује компатибилне залихе. Извршење би било успешно само ако би сваки посебно овлашћен хоп остао у својим границама цитата, границама, таксима, инвентаризацијом и политиком. Резултати сувапа нису доказали да је било који емитент касније испунио своје обавезе о ваучеру.

\subsection*{5.2 Предложене услуге маршрутизације и ребалансирања}
\addcontentsline{toc}{subsection}{5.2 Предложене услуге маршрутизације и ребалансирања}

Будућа услуга маршрута могла би подржати две различите активности.

\textbf{Извршење које је започео учесник.} С обзиром на улазне и излазне средства, износ и ограничења корисника, услуга би могла идентификовати пут и припремити извршење. Сваки хоп би имао свој одговорни Тржница, цитат, овлашћење, таксе, границе, инвентар и потврду. Атомске партије, ХТЛЦ-ови и гаранције су могући будући избор извршења, а не тренутно понашање протокола.

\textbf{Ребалансирање опционог Тржница.} Управници Тржница би могао објављивати циљеве за инвентар, дозвољене контрагенте, класе активи, границе за одклоњење котирања и границе по периоду. Одговорна служба могла би тражити компатибилне циклусе или ланце и извршити само овлашћене намере.

Ребалансирање би било оптимално. Покупни фонд може омогућити маршруте учесника док одбија ребалансирање излаза, или може омогућити само изабране средства, контрагенте и износе. Сваки извршен скок би изазвао квитанс, а било која такса за услугу би била објављена одвојено од такса за резервоар и протокол.

\subsubsection*{5.2.1 Конфедерација и оперативна способност}

Независни распоређивали би могли да управљају сопственим регистровима, интерфесема, услугама маршрута и политичким профилима, а истовремено би изабрали компатибилне стандарде података и примања. Извршавање преко профила остало би зависно од распореда.

Компатибилан профил би:

\begin{itemize}[leftmargin=*]
\item идентификују своје корене у регистру, операторе услуга, контролере и примените услове;
\item Публикују дозвољене и одбивене контрагенте, средства, адаптери и маршруте;
\item да примењују овлашћења, границе, таксе и ограничења за залихе свих учесника;
\item чувају доказе од цитата до примања по хопу; и
\item дозвољавају другим функционалним пулама да напусте или изабере другу регистру без брисања баланса или обавеза емитента.
\end{itemize}
Компатибилност може повећати доступне путање размене и смањити зависност од једног регистра или оператора. То не чини мрежу, CLC App, GEF или другу групу одговорним за испуњење емитента.

\subsection*{5.3 Предложен модел за мрежу и плату за услуге}
\addcontentsline{toc}{subsection}{5.3 Предложен модел за мрежу и плату за услуге}

Тренутни Protocol v1.1.0 наплаћује таксу по пулу и, када се конфигурише, додатну таксу по протоколу за директни суап по пулу. Ове текуће таксе остају различите.

Будући програм може одвојено добити:

\begin{enumerate}[leftmargin=*]
\item \textbf{мрежни удео}, дефинисан као утврђени део прикупљених накнада Тржница које учествују; и
\item у) \textbf{плата за маршрутизацију или услугу}, поплаћен за идентификован будући сервис.
\end{enumerate}
Предложена мрежа не представља додатни проценат који се примењује на целокупну суму суапа након што се већ рачуна такса по групи. За резервоар \texttt{p}:

\[
\tau_p = f_p \cdot r_p
\]

где је \texttt{f\_p} стопа покупне плате, а \texttt{r\_p} је предложен део те покупне плате која је додељена мрежном програму.

За одређени период:

\begin{itemize}[leftmargin=*]
\item \textbf{брутни трошкови за резервоар} су сума стварних прикупљених таксова за сваки фонд;
\item \textbf{Квитанте за мрежу} су пријављени дела тих прикупљених таксама;
\item \textbf{Квитанте за таксу за услуге} сакупљају се одвојено трошкове за маршрутизацију или услуге; и
\item \textbf{Квитанте плате за програм} једнаке квитанте за мрежу и квитанте за услуге.
\end{itemize}
Никакве категорије се не броје два пута. Текуће таксе протокола нису укључене, осим ако одвојену политику законски пренаправити стварне пријемке протокола таксе у будући програм.

За агрегиран приближавање, нека:

\begin{itemize}[leftmargin=*]
\item \texttt{Q\_swap} је вредност извршених суапова по групи за одређену кохорту и период; и
\item \texttt{$\tau$} је ефикасна стопа предложеног мрежног ракета и одвојено идентификоване таксе за услуге над извршеним вредностима суапа.
\end{itemize}
Онда:

\[
F \approx \tau \cdot Q_{\mathrm{swap}}
\]

Ово је аналитичка апроксимација, а не обећање прихода. Сваки улазак захтева одређену кохорту, период, јединицу, временски штамп процене, искључења и политику исправљања.

\subsubsection*{5.3.1 Квитанти који су прихватљиви у готовини и у природи}

Накнаде се могу приступити у средств. подстицајућим фунгибилним средствама или у ваучерима и другим материјалним средствама. Не може се аутоматски исплатити новчани трошкови или захтеви покривања. Свака размена или конверзија би захтевала овлашћење, доступни инвентар, откривене локације, границе и стварно извршење.

Нека је \texttt{$\chi$} реализовани део прихода плате који је прихватљив у готовини након ограничења политике, неуспелих конверзија и проласка. Квитанти који се могу користити у броју су:

\[
F_{\mathrm{cash}} \approx \chi \cdot F
\]

Бюджетна анализа и анализа равнотеже би користили реализовану \texttt{F\_cash}, а не бруто котиране таксе или номиналну вредност инвентаризације у природи. Будући програм ће одвојено износити бруто пошлине по пулу, приходе за мрежу, пошлине за услуге, састав имовине, резултате конверзије и приходе за коришћење новца.

\subsection*{5.4 Предложени програми ликвидности}
\addcontentsline{toc}{subsection}{5.4 Предложени програми ликвидности}

Будући, одвојено документовани програм ликвидности могао би да додељује средства одређеним пуловима или услугама рутинга. Тренутни \texttt{SwapPool} договори не штампају акције Пула или аутоматски стварају права на повраћање, повраћање, награду, управљање или добитак.

Сваки програм би објавио:

\begin{itemize}[leftmargin=*]
\item одговорни субјект и учествујући Управници Тржница;
\item донатовани средства и да ли је трансфер исплаћени, повлачиви, донирани или награђени;
\item уговор о чувању и техничком контролу;
\item дозвољене употребе, границе, блокирања, врата за повлачење и додељавање губитака;
\item прихватљивост таксе или подстицања и да ли је износ може бити нула;
\item пријаве, сукоби, жалбе и средства за опоравак; и
\item Миграција, прекид и третман преосталих илага и обавеза.
\end{itemize}
Материјални ризици укључују инвентар који је тешко разменити или испунити, ниску финансијску прихватљивост, неисправност емитента, неуспех у уговору или пружању услуга, промене управљања и ограничења излаза. Лимити, резерви, квитанти и ташкине могу смањити или открити неке ризике; Они не елиминишу губитак.

За екс-пост аналитичку метрику, нека:

\begin{itemize}[leftmargin=*]
\item \texttt{$\phi$} је реализовани део прихода плате за програм који су додељени по условним условима програма; и
\item \texttt{K} је мерена вредност актива који се покривају програмом по једној наведеној методи.
\end{itemize}
Онда:

\[
\mathrm{FeeFlow}_{LP} \approx \frac{\phi \cdot F}{K} = \frac{\phi \cdot \tau \cdot Q_{\mathrm{swap}}}{K}
\]

Ова метрика описује реализовани поток плате по мерељеном програмском активу. То није APY, прогноза, дивиденда или гарантована повратка. У извештајима би се одржавали одвојени обем суапа, презентација откупа, испуњавање емитента, ослобођење, трајање задржавања, губици, повлачења и квитансе плате.


\section*{6. Предложена ликвидност на нивоу мреже и управљање}
\addcontentsline{toc}{section}{6. Предложена ликвидност на нивоу мреже и управљање}

Опит из историјске Sarafu Network мотивисао је истраживање два шире питања дизајна:

\begin{enumerate}[leftmargin=*]
\item како независно регулисане пуле могу координирати услуге ликвидности и маршрутизације; и
\item како могу заједничке услуге да остану одговорне док учешће расте.
\end{enumerate}
Један предложени модел CLC-а уводи \textbf{предложену CLC Мрежну тржницу} за поравнање и усклађивање буџета на нивоу мреже и \textbf{предложени CLC токен управљања}. Ни једно ни друго не постоји у актуелној апликацији или Protocol v1.1.0. Друге примене могу користити задруге, јавне установе, федерације, multisig структуре, пружаоце услуга или друге одговорне структуре без усвајања било ког предлога.

\subsection*{6.1 Контроле и опционалност}
\addcontentsline{toc}{subsection}{6.1 Контроле и опционалност}

Тренутни Protocol v1.1.0 може применити ограничења баланса токена и контроле залиха. Ове контроле ограничавају одабране уговорне акције; Они не спречавају неуспех емитента, губитак вредности, губитак кључа, неуспех у договору или сваки други губитак.

У будућем распореду би могли да се додају прекидачи кола, временски блоки, резерви, гаранције или осигурање. За сваку заштиту би требало да буде идентификована одговорна страна, покривено догађај, финансирање, ограничење, искључења, трајање, доказ, процес захтева и применљиви услови.

Предложено управљање би требало да одржава ограничене и ревидиране овлашћења за регистар и службе. Обичне акције могу користити извештај, одлагање, објавување разлога, апелацију и исправне процесе. Спешне акције би требале да производе податке о инцидентима и да истекају или да добију ревизију у складу са усвојеном политиком.

Копатибилан пут може омогућити више размена између Тржница. Извршавање мулти-хопа и мерење захтевају додатни софтвер, овлашћења, границе, управљање неуспеха и управљање. Више цитираних или извршених суапа сама по себи не би доказала испуњење ваучера или социјални утицај.


\section*{7. Предложено CLC средства за управљање и управљање}
\addcontentsline{toc}{section}{7. Предложено CLC средства за управљање и управљање}

У овом поглављу је представљен један изборни дизајн управљања и координације ликвидности у будућности. Предлог CLC управљачког токена, stCLC, sCLC, управљачког третма, осигураног фонда, Водопада, мерника, мандата, прозора такс-кредита и програма спољног тржишта нису део тренутног CLC App или Protocol v1.1.0. Свака од њих би захтевала имплементацију, одговорног оператора, објављену политику и услове, ревизију безбедности и одговарајуће законско одобрење.

CPP не захтева овај модел токена. Компатибилни распоређивачи би уместо тога могли користити кооперације, јавне агенције, федерације, мултисигсе, услуге или друге структуре одговорне.

\subsection*{7.1 Циљ}
\addcontentsline{toc}{subsection}{7.1 Циљ}

Ако се усвоји, предложен слој управљања би могао:

\begin{itemize}[leftmargin=*]
\item координира заједничке регистре и услуге маршрута;
\item додељује одобрете мандати за ликвидност на независно управљане пулове;
\item управља изборним, изричито одређеним, финансираним програмом покривања;
\item одржавање заједничке инфраструктуре и посматраљивости;
\item признају доприносиоце према објављеним правилима о прихватљивости и борби против игара; и
\item одржавање прегледа, апелације, транспарентности и кредибилног излаза док учешће расте.
\end{itemize}
\subsection*{7.2 Предложен модел управљања и средства}
\addcontentsline{toc}{subsection}{7.2 Предложен модел управљања и средства}

Дизајн садржи три различите предложене средства:

\begin{enumerate}[leftmargin=*]
\item \textbf{Предложена CLC токен управљања:} трансферни базови управљачки активи на основу усвојена правила емисије, чувања, гласања и усклађивања.
\item \textbf{stCLC:} непремењива квитанца за гласање, која се чека када се блокира предложен токен управљања CLC. Она би представљала власт управљања ограничена временом и могла би омогућити откривање делегације.
\item \textbf{sCLC:} укупни активи који су издадени у оквиру политике. Она би могла подржати ограничени приступ такси-кредиту или одобрена ликвидност и подстицања за маршрутизацију и могла би истећи или бити спаљена на крају периода.
\end{enumerate}
Ни stCLC ни sCLC не би били власнички капитал, дивиденд, остатак захтева, обећање повратака или заједнички ваучер. Прихваћена политика могла би да постави емисију sCLC и приступ такси-кредиту на нула у било коју епоху.

Пре употребе би се објављивали блокирање управљања, минимални трајања, смањење температуре, делегације, ограничења концентрације и прагови критичних акција. Разблоковани предложени токенци CLC управљања не би гласали по овом дизајну.

\subsubsection*{7.2.1 Илустративна понуда и додељавање}

Илустративна укупна количина CLC токена управљања је \textbf{500,000,000}. То је параметар модела, а не актуелно издавање, понуда, вредновање нити обећање ликвидности.

Илустративна додеља је:

\begin{itemize}[leftmargin=*]
\item \textbf{15\% --- Grassroots Economics Foundation:} трајно стављено, непремењиво и везано за идентификован GEF мултисиг према објављеним правилима потписивања и ротације.
\item \textbf{15\% --- основни тим и рани партнери:} уложили су се током одређеног политичког клифа и линеарног периода улажења од 24 месеца, са унапред објављеним правилама гласања и преноса.
\item \textbf{30\% --- приватне награде:} Признање за ране доприносиоце који пружају одобрену мрежну ликвидност, са одређеним политичким правомоћом од 24 месеца.
\item \textbf{40\% --- резерва јавне ликвидности:} одржава се у управном третма у времену за изборне програме спољног тржишта и доступности.
\end{itemize}
У оквиру илустративних контрола јавне ликвидности:

\begin{itemize}[leftmargin=*]
\item не би више од 10\% укупне снабдевања било активно на спољним местима истовремено;
\item сваки распоређивање би истекао након 90 дана, осим ако се не понови кроз усвојен процес;
\item позиције ликвидности и било који токен позиције би остали под контролом откривеног управљања; и
\item предложени токенси CLC управљања који се користе у позицијама ликвидности не би гласали осим ако нису блокирани по истим правилима као и други токенси за гласање.
\end{itemize}
Будући лансирање могло би почети са откривеном мултисигом и прелазом само након одвојеног одобрења мелиона за ојачавање, као што су независни ревизије, праћење, путнике инцидента и тестиране процедуре паузе и форка. Сваки прихваћен параметр и процес промене би били објављени одвојено.

\subsubsection*{7.2.2 Стадије доступности и додељавања}

У будућем програму давања мог би се користити фазиран прозор доприноса и објављена референтна вредност за пријем и буџет. Та референција не би била обећање тржишне цене, процене, повратака или напуштања.

У условима доприноса би се навело да ли су имовине дониране, дарене, извучене или враћене; који их контролише; њихова дозвољена употреба; затварачи; распоређивање губитака; извештавање; и условима излаза. Велики доприноси могу се суочити са ограничењима гласања или смањењем тежине управљања.

Свака ликвидност на спољном тржишту би захтевала посебну одлуку која би идентификовала локације, одговорне операторе, средства, ограничења, распореда, конфликте, контроле ризика, извештавање и прекид. Политика не би намењена ни гарантовала цену.

\subsubsection*{7.2.3 Предложена програма за посев удара}

Будући програм би могао да додељује део третма управљања доприносницима који стављају средства у одређену Тржнице обавеза и мериво побољшавају приступ берзи и испуњавање потврђеног емитента.

Прихваћена политика прихватљивости би:

\begin{enumerate}[leftmargin=*]
\item идентификују одобрене пулове, имовине, минимални трајање и услови за закључавање;
\item дефинише продуктивни допринос користећи доказе извршене суапе и посебно доказану презентацију, испуњење и ослобођење;
\item мере маргинални допринос, а не укупну вредност која је искључиво блокирана;
\item искључују самопродајућу и кружну активност прања;
\item примјењују границе за сваки ентитет и смањујућу доходност; и
\item обезбедити прозор посматрања, спора и исправљања пре коначног додељавања.
\end{enumerate}
Сваки грант остаће предмет објављених услова програма, усвајања, допустимости, испуњавања и правила губитка.

\subsection*{7.3 Предложено гласање, подстицања и приступ таксама}
\addcontentsline{toc}{subsection}{7.3 Предложено гласање, подстицања и приступ таксама}

Према овом дизајну, власници stCLC могли су да гласају за прихватљиве пулове, мандати маршрутних услуга, подељени буџети или друге усвојене предлоге. Система курисаних мерка могла би превести гласове у ограничени sCLC подстицаји за прихватљиве ликвидности или операторе маршрутизације.

Политике мерења одржавају извршене суапе и испуњавање емитента одвојене. Она не би наградила цитиране, већ неиспостављене путеве, нерешене презентације, самопродавање или укупну вредност затворена без доказа корисне активности.

Прихваћен процес управљања такође би могао објављивати епохијски буџет за плаћање кредита. Добиљни власници stCLC могу добити sCLC дозвољавање ограничених, временски ограничених суапа од одређених пула за држање таксама у допуњене имовине или програме. То би био приступ доступним инвентарским подацима, а не пасивни приход или власништво резерва за придржавање такса.

\subsubsection*{7.3.1 Извршавање таксе за услуге и избор профила}

У будућем регистарском профилу може бити потребно да учествујуће пулове или службе маршрута користе компатибилан адаптер за плату или крпка. Профил може одбити откриће, маршрутизацију или подршку програму када захтеван потврда не доказује прикупљање плате.

Имена као што су \textbf{PoolFactory } или \textbf{ ФиХук} описати могуће будуће модуле; Они нису Protocol v1.1.0 договори. Свака имплементација би открила код, адресе, контролере, примаоце, стопе, прихватљиве трансакције, обраду неуспеха и статус ревизије.

Извршење би било специфично за профил. Полу који је искључен по једном профилу могао би наставити да функционише локално или да се придружи другом компатибилном регистру ако су његови договори и зависности остали функционални. Клијенти би требали да открију доступне профиле и да покажу одговарајуће трошкове пре одобрења.

\subsubsection*{7.3.2 sCLC буџет и изборне контроле спољашњег плавања}

Након што се финансирање усвоји бюджети са већим приоритетима, у будућем процесу би се могао објављивати епохијски буџет за такс-кредит \texttt{F\_epoch}, укључујући нула. Политике могу да додељују лимит учесника пропорционалан гласачкој моћи stCLC:

\[
\mathrm{limit}_{\mathrm{user,epoch}} = F_{\mathrm{epoch}} \times \frac{\mathrm{stCLC}_{\mathrm{user}}}{\mathrm{stCLC}_{\mathrm{total}}}
\]

Свако окно би остало подложено разрешавањима, ограничењима, инвентаризацијом, прихватљивошћу, правилима инцидента и примењивим законима.

Одвојна политика би могла да одобри ограничене, временско тежине прикупљање предложеног токена CLC за управљање само како би се смањила мерена површина напада спољног управљања. То би захтевало:

\begin{itemize}[leftmargin=*]
\item објављен триггер заснован на спољном плутању током одређеног периода;
\item максимални део реализованих, допустивих таксама и обема места;
\item Извршавање уз временску тежину без прецизног обавештења о времену;
\item аутоматски заустављање када се не испуњавају усвојена покривеност или прагови за рад;
\item поношење или стављање стечених токова на откривен негласни рачун; и
\item експлицитна изјава да програм не циља или не гарантује цену и не утиче на испуњење емитента.
\end{itemize}
Полиција би могла остати онеспособљена на неопредељено време.

\subsubsection*{7.3.3 Портфолио Пулс и потврде}

А \textbf{Портфолио Паул} Била би обична Тржница обавеза курирана око мисије или службеног домена. Његов Управник Тржнице може бити појединца, кооперација, заједничка група, јавна агенција, федерација, мултисиг, оператор услуга или друга структура одговорна.

Поддршка се може добити путем:

\begin{enumerate}[leftmargin=*]
\item директне доприносе у складу са објављеним одредбама фонда;
\item временски ограничени мандати за ликвидност одобрени кроз усвојен процес управљања; или
\item sCLC усмерен приступ ограниченим буџету након водопада када је тај механизам омогућен.
\end{enumerate}
Портфолио пули би остали независни и могли би користити заједничке или независне регистре.

Сертификати би могли да информишу пријем или третман ризика, али не би створили право на таксе, профит или остатак имовине. Размештај може користити:

\begin{itemize}[leftmargin=*]
\item \textbf{потврде везане за регистар} од идентификованог верификатора који користи објављену методу; или
\item \textbf{купонови за проверну услугу} које представљају уговор који се може изкупити за обављање пријављене ревизије или проценке.
\end{itemize}
Пул би открио како потврда утиче на пријем, методе размене, границе или допустимост и како се решавају грешке, истека времена, сукоби и жалбе.

\subsection*{7.4 Предложени водопад и буџети}
\addcontentsline{toc}{subsection}{7.4 Предложени водопад и буџети}

Предложени водопад ће одредити само реализоване, прихватљиве квитанте у складу са усвојеном процесом управљања. Она би разликовала:

\begin{itemize}[leftmargin=*]
\item \textbf{Квалификовани средства (\texttt{E\_cash}):} Квалификовани фунгибилни активи који се могу користити или конвердовати у складу са политиком за облигације у готовности; и
\item \textbf{Укупни средства (\texttt{E\_kind}):} купонови или други активи који могу подржати борсу у мрежи, али се не рачунају на покривеност у готовини или оперативне обавезе, осим ако се заправо конвертују.
\end{itemize}
Политика конверзије би идентификовала овлашћене операторе, имовине, локације, изворе цена, временске прозоре, границе проклањања, ограничења, конфликте, записи и процедуре инцидента. Конверзија казначейства не би одређивала обавезу емитента о ваучеру или методу разменног курса групе.

Илустративни поредак приоритета је:

\begin{enumerate}[leftmargin=*]
\item \textbf{Циљ финансираног покривања:} изградити сваку усвојену резерву или предложену осигуравајућу фонду према свом откривеном циљу за дефиниране покривене догађаје.
\item \textbf{Основне операције:} финансирање ограниченог, откривеног буџета за правни рад, образовање, комуникације, инфраструктуру, ревизије и посматраност.
\item \textbf{Мандати за ликвидност:} додељују одобрена средства за одређене сврхе пула или службе маршрута у оквиру прегледе ограничења и заласка сунца.
\item \textbf{Опционална контрола спољашњег пливања:} финансирање само када су испуњени објављени триггерски и приоритетни прагови; у супротном додељују нуле.
\item \textbf{Предложен буџет за такс-кредит:} додељују било који одобрен остатак одређеним пулима за држање таксама и објављују \texttt{F\_epoch}, који може бити нула.
\end{enumerate}
Одлуке о буџету могу узети у обзир потврђено испуњење, покривено изложеност, допустиве резерве, ограничене употребе, извршене маршруте, концентрацију, инциденти и резултате гаранта. Свака метрика би следила правила доказа и кохорте из Додатка Ц.

Могући поток доприносника би био:

\begin{enumerate}[leftmargin=*]
\item доприносници преносе прихватљиве средства по одвојеног одобрењем статута или условима програма;
\item прихватљиви учесници добијају и, ако је то дозвољено, закључавају предложене токове управљања CLC како би добили stCLC;
\item реализовани прихватљиви квитанти за раке или таксу за услугу улазе у водопад;
\item фондови водопада усвојили су приоритете у редоследу; и
\item само активирана политика након Ватерфала издаје sCLC или отвара ограничено окно кредитовања таксама.
\end{enumerate}
Ниједан корак не ствара права власништва, повратака, повлачења, покривености, награде или управљања изван оних што су експлицитно наведене у примењујућим условима.


\section*{8. Технички опсег и раст}
\addcontentsline{toc}{section}{8. Технички опсег и раст}

У овом поглављу се описује изборни или предложен посао. То није листа карактеристика које се гарантују да су присутне у току CLC App или Protocol v1.1.0.

Могуће области рада укључују:

\begin{itemize}[leftmargin=*]
\item извршење рутовања преко компатибилних пулова, са регистрацијом, цитирањем, лимитом, таксом и откривањем инвентара;
\item Уколико се не може постићи атомско расплатавање, прилагођавачи за криздоменско извршење који су у времену блокирани или HTLC;
\item интерфејсе и алате политике за мале или личне базе;
\item ревидирани регистри по ваучерима, пулима, методама размене, лимитима, контролерима и таксама;
\item Контактори пружаоца плаћања специфични за распоређивање, текови чекања и контроле прихватљивости;
\item повезаности фунгибилних актива са спољним локацијама ликвидности за ребалансирање и ликвидност плаћања; и
\item конверзија казначейских средстава која се покрива политиком за усвојена покривеност, оперативне трошкове или мандати за ликвидност.
\end{itemize}
Цене на спољном тржишту не би одређивале колико је емитент дужен по условима ваучера. Поул би могао користити заштићену спољну референцију за фунгибилан актив, али његова објављена метода размене, границе, таксе и инвентар ће управљати његовим котисацијама.

\subsection*{8.1 Предложени стандарди за маршрутне услуге и SDK}
\addcontentsline{toc}{subsection}{8.1 Предложени стандарди за маршрутне услуге и SDK}

\textbf{Откривање.} Предложена услуга маршрута би тражила идентификоване регистре за пријем имовине, методе размене стопе, границе, таксе, инвентар, инциденти и информације контролера. Кашеви записи укључују границе свежести и идентификатори извора.

\textbf{Профили мреже.} Клијент може подржати више од једног корена регистра или профила политике. Она би учеснику рекла који профил, контрагенте, адаптери, ограничења и одговорни оператори услуга користи цитат. Кружпрофилни маршрут би морао задовољити све важеће услове за хоп.

\textbf{Политика пута.} Одговорни оператор може искључити несигурне зависности или контрагенте и применити границе на нивоу трасе, захтеве за свежину и здравствене критеријуме. Ови сигнали би подржали одлуку; да не гарантују испуњење или заштиту од губитка.

\textbf{Таксе и границе.} Цитат би изложео таксе за Поул, било које додатне текуће таксе за Протокол и било које одвојено предложене таксе за маршрутизацију или услуге. Извршење би одбацило истекале цитати или кршење граница.

\textbf{Атомност и опоравак.} Мулти-хоп екзекуција би била атомска где је то могуће. Уколико је користила ХТЛЦ-а или ескрове, служба би откривала временске излазке, траке прекида, одговорне контролере, процедуре инцидента и преостале ризике.

\textbf{Предложена сетовања и ребалансирање.} Услуга оптималног уласка може прикупљати намере на ребалансирање и тражити компатибилне циклусе или ланце. То би:

\begin{enumerate}[leftmargin=*]
\item Публикују машински читати квитанс који идентификује извршене циклусе, средства, износе, временске марке за процене и таксе;
\item извршава усвојена ограничења за период и политику контрагенте;
\item одбацивање активности које крше овлашћења, ограничења или доступни инвентар учесника; и
\item чувају детерминистичке уносе и поступке за преиспитивање и решавање спорова.
\end{enumerate}
\textbf{Захтеви за SDK.} SDK за извршене маршруте би обезбедио детерминистичко мапирање цитата до примања, контроле инваријанта по хопу, разумне кодове неуспеха и регистроване регистрове за ревизију. У току Protocol v1.1.0 \texttt{SwapRouter} се налазе само цитати; није извршила ове предложене маршруте.

\subsubsection*{8.1.1 Минимална спецификација компатибилности конфедерације}

Екосистем "Полл" који тражи маршрутизацију преко профила би објавио информацију која се може читати машином за:

\begin{enumerate}[leftmargin=*]
\item \textbf{Корени регистра:} Идентификатори за имовине, пулове, методе размене, границе и политике такса, или један корен који их детерминистички решава.
\item \textbf{Квитанце:} Профил, активи у и из, количине, извор цитата и временски штамп, гранични снимк, таксе, резултат инвентаризације и резултат извршења за сваки скок.
\item \textbf{Оперативни сигнали:} свежи ограничени подаци о инвентаризацији, ограниченој употреби, инцидентима и било ком посебно доказаном испуњавању или финансираном заштити.
\item \textbf{Политички ограничења:} дозвољене или одбивене контрагенте, класе активија, адаптери и било који захтев за гаранцију.
\item \textbf{Кодови за неисправност:} детерминистичке објашњење одбијања, истека, ограничења, инвентаризације, политике, зависности или неуспеха инцидента.
\end{enumerate}
Профил може додати покривеност, поштовање, арбитраж или друге услуге без тога да их чине захтевима за основну CPP компатибилност. Свака изборна служба би идентификовала одговорну особу, овлашћење, опсег и услове.

\subsection*{8.2 Лицензирање, верификација и излазак}
\addcontentsline{toc}{subsection}{8.2 Лицензирање, верификација и излазак}

Protocol v1.1.0 договори су EVM-Спостављени. Уговора у каталогу \texttt{src} регистар Протокола објављују се под именом AGPL-3.0, осим за идентификоване непроменене компоненте треће стране које задржавају своје услове. Публиковани извор, АБИ и инструкције за распоређивање подржавају независни преглед, али сами по себи не доказују ревизију, сигурно распоређивање или правно поштовање.

Свака дислокација ће одвојено откривати своју верзију кода, порекло изграде, адресе, овлашћења контролера и ажурирања, статус ревизије, огледала регистра и било које заштите од временског блокирања или паузе.

Предложена \textbf{комплект за вилицу} може укључивати:

\begin{enumerate}[leftmargin=*]
\item детерминистичке скрипте за распоређивање;
\item регистарске снимке и алате за извоз;
\item Документовани процес за повлачење услуга маршрута, SDK-а и интерфејса на нови корен регистъра;
\item контролна листа Управник Тржнице за безбедно напуштање заједничког регистра; и
\item Миграциону контролну листу за неизплаћене ваучере, укључујући обавештења емитента, срокове представљања и испуњавања, континуиран приступ записима и средства за опоравак.
\end{enumerate}
Компатибилни вилки могу побољшати отпорност када заједници, кооперативи, јавне агенције, федерације, мултисиг или оператори услуга требају другачије управљање. Стварна континуитет још увек зависи од власништва уговора, кључева, зависности, интерфејса, инфраструктуре, правних обавеза и услуга трећих страна.


\section*{9. Предложена економија за програме ликвидности}
\addcontentsline{toc}{section}{9. Предложена економија за програме ликвидности}

У овом поглављу се описује будући модел који је усвојен одвојено. То није тренутна функција апликације, понуда, обећано повратак или право креирано депозирањем у \texttt{SwapPool}.

\subsection*{9.1 Предложени извори прихода}
\addcontentsline{toc}{subsection}{9.1 Предложени извори прихода}

Будући буџет за мрежу може добити:

\begin{enumerate}[leftmargin=*]
\item откривено \textbf{мрежни рајк} узимају се као диел укупљених таксама учесничких пула;
\item одвојене трошкове за маршрутизацију или услуге од реализованих заједничких услуга; и
\item други изричито усвојити, примљени приходи.
\end{enumerate}
Укупне трошкове резерви које задржавају резерви не представљају приход из мрежа. Тренутни Protocol v1.1.0 користи другачији модел: опционална такса протокола се додаје такси групи и директно се šalје конфигурисаном примаоцу.

\subsection*{9.2 Илустративна раке математика}
\addcontentsline{toc}{subsection}{9.2 Илустративна раке математика}

Ако учествујућа група покупља 2.00\% покупну таксу, а усвојена мрежа покупља 20\% ове покупне таксе, предложена ефикасна стопа покупне таксе за мрежу на маршрутизованој вредности је:

\texttt{rake\_rate = 2.00\% $\times$ 20\% = 0.40\% = 40 bps}

Ако је 25\% примјењених раке и средстава за услуге допустимо и конвертибилно за пријављену употребу у новчаној номинацији након трошкова, удео у новчаној употреби раке од 40 бепса је око 10 бепса.

Предложена приход мрежа је:

\texttt{network\_rake\_received + routing\_or\_service\_fees\_received}

Не додајте бруто пошлине за пуле у мрежну раке: раке је пренос од тих пошлина и иначе би се рачунао два пута.

\subsection*{9.3 Право на програм ликвидности}
\addcontentsline{toc}{subsection}{9.3 Право на програм ликвидности}

Одвојено прихваћен програм може финансирати инвентаризацију пула, услуге маршрутизације, праћење или друге мандати. Уредби би морали да открију:

\begin{itemize}[leftmargin=*]
\item да ли је трансфер дар, донација, позајма, принос који се може вратити или куповина;
\item чување и контролу;
\item правила повлачења, повлачења, губитка и приоритета;
\item прихватљивост плате и награде;
\item права на управљање;
\item методе процене и извештавања; и
\item суспензија, прекид и средства за опоравак.
\end{itemize}
Садашњи \texttt{SwapPool} не ствара токен за Тржница-share или аутоматско право на допринос. Сваки пост-метрички показаци морају бити засновани на оствареним приходима и губицима, не морају бити представљени као обећани доход и могу бити нулеви или негативни.


\section*{10. Свеобухватни оквир за ризике}
\addcontentsline{toc}{section}{10. Свеобухватни оквир за ризике}

Овај предложен оквир дели десет категорија ризика. За сваку категорију она идентификује могуће индикаторе, контроле, тестове стреса и индикативни апетит за ризиком. Ово су препоруке за дизајн, а не тврдње да је свака контрола распоређена, ефикасна или довољна. Лимити, резерви, гаранције, мониторинг, покривеност и управљање не могу елиминисати губитак.

\subsection*{10.1 Протокол и ризик од паметних уговора}
\addcontentsline{toc}{subsection}{10.1 Протокол и ризик од паметних уговора}

\begin{itemize}[leftmargin=*]
\item \textbf{Загроза:} контрактне грешке, грешке у унапређењу, неуспехе у зависности и погрешно конфигуриране границе или накнаде.
\item \textbf{Индикатори:} Резултати ревизије, необјашњиви покрети инвентара, неизменни неуспехи и необични повраћаји.
\item \textbf{Могуће контроле:} независне ревизије; минималне привилегијске улоге; оповестени контролери за прокси и зависности; модернизације у временском времену; мониторинг на ланцу; паузе инцидента са објављеним овлашћењем и критеријумима; и тестиран пут миграције.
\item \textbf{Испитивање стреса:} недоступне зависности од котирања или ограничења, недостатак инвентара, паузе у уговорима и експлозиви саобраћаја.
\item \textbf{Апетит за ризиком:} минимално пре повећања неизплаћених обавеза или облика суапа.
\end{itemize}
\subsection*{10.2 Економски и тржишни ризик}
\addcontentsline{toc}{subsection}{10.2 Економски и тржишни ризик}

\begin{itemize}[leftmargin=*]
\item \textbf{Загроза:} Тешки инвентар, једнострани потоци, брзи повлачење и манипулисани или устарели референци цена.
\item \textbf{Индикатори:} Висока употреба баланса токена, ширење разлика у котисању, често одбијање границе и концентрисан инвентар.
\item \textbf{Могуће контроле:} тренутне границе баланса токена по пулу; предложени ограничења за ролинг или рачуна; резерве које су одвојене одвојено; чуване референце за цене; искључења маршрута; и временски ограниченим таксима за инциденти или ограниченима.
\item \textbf{Испитивање стреса:} велики кретања референтних цена, повећања презентације, захтеви за повлачење и прекидове извора података.
\item \textbf{Апетит за ризиком:} умереност само у оквиру објављених параметара и финансираног капацитета за износ губитака.
\end{itemize}
\subsection*{10.3 Ризик емитента и купона}
\addcontentsline{toc}{subsection}{10.3 Ризик емитента и купона}

\begin{itemize}[leftmargin=*]
\item \textbf{Загроза:} емисија изван капацитета испуњавања, неизплата емитента, погрешни услови и слабо одређени прозори доступности или презентације.
\item \textbf{Индикатори:} смањења стопа потврђеног испуњавања, старење неизплаћених купонова, нерешене жалбе и концентрисана експозиција према једном емитенту.
\item \textbf{Могуће контроле:} условна пажња емитента; јасни ваучер и услове понуде; границе емисије или допуштања; независно финансиране облигације или гаранције; извештавање кохорти; и одговорног прегледа регистра.
\item \textbf{Испитивање стреса:} Имитентска неплаћност, регионални шокови производње, фалсификовани захтеви и дуготрајни одлази у испуњавању.
\item \textbf{Апетит за ризиком:} Пониже се са повећањем концентрације емитента или понуде.
\end{itemize}
\subsection*{10.4 Ризик представљања изкупа и испуњавања}
\addcontentsline{toc}{subsection}{10.4 Ризик представљања изкупа и испуњавања}

\begin{itemize}[leftmargin=*]
\item \textbf{Загроза:} неважни или дуплирани подаци, недостатак капацитета емитента, резерви, недостатак логистике и недостатак података о ослобођењу.
\item \textbf{Индикатори:} Латенција у испуњавању захтева, неуспешни или спорни подаци, резервирање, западни бројеви билета и испуњене јединице које нису испуњене.
\item \textbf{Могуће контроле:} објављене процедуре представљања и испуњавања; откривање капацитета; више места испуњавања, ако је то законско; стандарди доказа; путем пријаве и исправљања; и записи који разликују презентацију, испуњавање и ослобођење.
\item \textbf{Испитивање стреса:} 2 до 4 пута већи обем презентације, прекид објеката, неуспех снабдевача и покушаји дуплирања употребе.
\item \textbf{Апетит за ризиком:} Мало где су укључени неопходни производи, ранљиви учесници или дуги прозори испуњавања.
\end{itemize}
\subsection*{10.5 Ризик управљања}
\addcontentsline{toc}{subsection}{10.5 Ризик управљања}

\begin{itemize}[leftmargin=*]
\item \textbf{Загроза:} Стјуард или контролер улазак, брзе промене параметара, сукоби интереса, скривене техничке моћи, и слабе апелације.
\item \textbf{Индикатори:} концентрисана власт, честе хитне акције, необјашњене промене политике и понављање повратних повлачења.
\item \textbf{Могуће контроле:} откривање улога и моћи; пропорционални прагови одобрења; временски часови; сукоби и правила одбијања; јавни записи промена; жалбе; аутоматске западне сунце за хитну енергију; и виочност.
\item \textbf{Испитивање стреса:} Пропозиције противречности, губитак потписница, покушаји подкупа и улазак путем акумулације или делегиране контроле.
\item \textbf{Апетит за ризиком:} мало за акције које утичу на методе вредности, повлачење, корен регистра, покривеност или овлашћења за хитну ситуацију.
\end{itemize}
За будуће распоређивање предложеног CLC управљачког токена, тест за улазак би укључио акумулацију, а затим покушаје пренаправљања буџета, ослабљења стандарда регистъра, одобрења мандата повезаних страна или исцрпљења финансираног покривености. Могуће гаранције укључују блокирање управљања, већи праг за критичне акције, одлагање спровођења, транспарентно праћење делегације и концентрације, процес инцидента и веродостојни поступак одласка. Ниједна није представљена као распоређена само појављујући се овде.

\subsection*{10.6 Правни и риски у складу са законом}
\addcontentsline{toc}{subsection}{10.6 Правни и риски у складу са законом}

\begin{itemize}[leftmargin=*]
\item \textbf{Загроза:} купон, услугу, промоцију или управљачки актив који добија неочекивано регулаторно третирање; неуспех у заштити потрошача; изложеност на прање новца или санкције; и трансграничних ограничења.
\item \textbf{Индикатори:} знаме јурисдикције, истраге регулатора, жалбе, утакмице ограничених страна и разлика између рекламираног и стварног понашања.
\item \textbf{Могуће контроле:} преглед по категорији и надлежности; точне откриће; геоосновани интерфејси; пропорционална проверка допустимости или потврде; контроле промоције; записи о овлашћењу и прихватању; и јасне одговорне стране.
\item \textbf{Испитивање стреса:} ограничење надлежности, прекид пружања услуга, обавезна прекласификација и налог за заустављање функције или класе актива.
\item \textbf{Апетит за ризиком:} ниски; ограничити или зауставити неподржну активност.
\end{itemize}
\subsubsection*{10.6.1 Правно позиционирање и третман предложених токова}

\begin{enumerate}[leftmargin=*]
\item \textbf{Проверљива инфраструктура:} Protocol v1.1.0 договори су EVM- компатибилни и њихов извор је објављен на основу лиценца и изузетка трећих лица утврђених у регистру протокола. Публикација омогућава ревизију, али сама по себи не доказује ревизију, сигурно распоређивање или правно поштовање. Свака дислокација треба да открије своју верзију кода, порекло изграде, адресе, овлашћења контролера и статус ревизије.
\item \textbf{Предложена позора знака:} у овом дизајну, предложен токен CLC управљања би координирао управљање и приступ политике. То не би створило дивиденде, поделу профита, остатке права или гаранцију вредности или ликвидности.
\item \textbf{Сврзани предложени активи:} одвојено имплементисана политика би могла омогућити блокирање предложеног токена CLC за управљање stCLC и могла би одобрити sCLC за епоху. У условним условима би требало да се прецизно дефинишу њихова права, ограничења, истека, преносимост и третман.
\item \textbf{Комуникације:} материјали за било који предложен CLC управљачки токен, stCLC или sCLC распоређивање не би требало да обећавају профит, процене, пасивни приход или гарантован приступ.
\item \textbf{Контроле надлежности:} за распоређивање може бити потребна геофенцирање интерфејса, потврде за ограничене класе, ограничења промоције, преглед специфичног актива или контроле који заустављају предложну функцију.
\item \textbf{Уведомљење учесника:} условни услови који би се усвојили би објаснили када се приступ може смањити или онемогућити по законским, оперативним или ризичним разлогама и да ли се примењује било која компензација или правна средства.
\end{enumerate}
\subsection*{10.7 Ризик маршрутизације и преко домена}
\addcontentsline{toc}{subsection}{10.7 Ризик маршрутизације и преко домена}

\begin{itemize}[leftmargin=*]
\item \textbf{Загроза:} делимично извршење, заглављени скочи, експлоатације мостове или гаранције, устареле котисације, инфлацију на путу и претходни рад најављених промена вредности.
\item \textbf{Индикатори:} стопе истека трајека, одстаци на депозитима, разлике у цитирању до извршења, понављање непотребних скока и инциденте у мосту.
\item \textbf{Могуће контроле:} атомско извршење, када је доступно; конзервативне HTLC или временске излазке из депозита; пут и политику контрагенте; границе на нивоу маршрута; детерминистичко картографирање котирања до примања; и одговорних оператора услуга.
\item \textbf{Испитивање стреса:} Спрема моста, реорганизација ланца, прекид зависности и један неуспешен скок на предложеној маршрути.
\item \textbf{Апетит за ризиком:} ниски до умерени само за идентификоване и контролисане зависности.
\end{itemize}
\subsection*{10.8 Ризик чувања и управљања кључевима}
\addcontentsline{toc}{subsection}{10.8 Ризик чувања и управљања кључевима}

\begin{itemize}[leftmargin=*]
\item \textbf{Загроза:} изгубљене или компрометисане кључеве, колузија потписника и нејасна власт за опоравак или администратор.
\item \textbf{Индикатори:} аномалне потписе, промене контролера, неуспешне ротације и необичне повлачења.
\item \textbf{Могуће контроле:} дозвола за више знакова или прага; кључеви подржани хардвером; раздвајање улога; ротација потписница; инвентаризације јавних контролера; контролисане границе повлачења; и тестиране процедуре за опоравак.
\item \textbf{Апетит за ризиком:} Мало је.
\end{itemize}
\subsection*{10.9 Репутација и социјални ризик}
\addcontentsline{toc}{subsection}{10.9 Репутација и социјални ризик}

\begin{itemize}[leftmargin=*]
\item \textbf{Загроза:} погрешне тврдње, штетне подстицаје, неприступне жалбе, лоше искуство испуњавања, пропуске приватности и заштите које фаворизују инсайдере.
\item \textbf{Индикатори:} жалбе по кохортима, нерешену спору, трендове у односу на перформансе емитента, концентрацију користи или губитака и повратне информације заједнице.
\item \textbf{Могуће контроле:} откривања на једноставном језику; начине за пријаву и исправљање; доступне доказе; транспарентно извештавање; преглед инцидента; и пропорционалне санкције за лажно представљање.
\item \textbf{Апетит за ризиком:} мало, са посебном бригом о погођеним заједницама и ранљивим учесницима.
\end{itemize}
\subsection*{10.10 Ризик концентрације и фрагментације}
\addcontentsline{toc}{subsection}{10.10 Ризик концентрације и фрагментације}

\begin{itemize}[leftmargin=*]
\item \textbf{Загроза:} зависност од малог броја емитента, пула, контролера, провајдера, мрежа или некомпатибилних вилица.
\item \textbf{Индикатори:} мере концентрације по емитенту, пулу, инвентарству, администратору или пружаоцу услуга; провалице пута између кластера; и критичне појединачне зависности.
\item \textbf{Могуће контроле:} објављени прагови концентрације; више одговорних оператора; компатибилни стандарди; независни регистри; тестиране процедуре излаза; и сигурну оперативну сарадњу.
\item \textbf{Испитивање стреса:} губитак највећег емитента, пула, оператора, пруживача или корена регистра.
\item \textbf{Апетит за ризиком:} специфичне за распоређивање и објављене, са строже ограничења за неопходне услуге или незамениве зависности.
\end{itemize}

\section*{11. Механизми управљања}
\addcontentsline{toc}{section}{11. Механизми управљања}

У овом поглављу се предлаже шаблон управљања. То не представља да се тренутни CLC App користи гласање о управљању токовима, временске блокирања, заједничко осигурање, процес захтева или свака контрола описана испод. Свако распоређивање би требало да идентификује своје стварне доносеље одлука, овлашћења, уговор, процесе и политике.

\begin{itemize}[leftmargin=*]
\item \textbf{Уставне вредности:} бригу о људима, бригу о животној средини, праведност, реципрочност, недоминирање и отпорност.
\item \textbf{Типови предлога:} промене таксе, лимита и индекса; мандати за ликвидност; Листирање и уклањање резервоара; одлуке о опционалном покривању; и ограде за параметре.
\item \textbf{Отговорни процес:} Употреба $\to$ процена $\to$ преглед ризика $\to$ одобрење $\to$ временски ограничење где је потребно $\to$ извршење. Оглашавање може доћи од управитеља, кооператива, јавних агенција, федерација, мултисиг, гласања на ланцу или друге откривене и одговорне структуре.
\item \textbf{Прагови одобрења:} параметризовани по акционим класима, са већим праговима за промене индекса вредности, овлашћења за хитне ситуације и друге критичне акције.
\item \textbf{Делегација:} изборна делегација са јавним мандатима, откривање конфликата и повлачење.
\item \textbf{Спремачи кола:} хитне паузе са одређеним критеријумима, овлашћеним операторима, условима настави и потребним пост-мортем.
\item \textbf{Пространост:} објављене промене и потоке, са посебним доказима за расплату суапа, испуњење емитента, резерве, ограничено коришћење, маршрутизацију и гаранте.
\end{itemize}
\textbf{Управљање регистром.} CPP- компатибилна дистрибуција може одржавати регистре открића за ваучере, токове и пулове. Ауторизоване контроле могу додати, ажурирати, суспендирати или уклонити регистарске уносе кроз раскључен процес управљања распоређивања. Узимање регистра утиче на откривање и маршрутизацију кроз тај регистар; не брише само по себи токен, не мења баланс власника, не испуњава обавезе емитента или не онемогућава уговор који иначе функционише.

Објављени правила регистра треба да обусловљају статус и могу идентификовати понављање неиспостављања, преваре или погрешна изјава, несигурно понашање у договору или трајно кршење објављених принципа као разлог за суспензију или уклањање. Уколико је то могуће, процес би требало да обезбеди обавештење, прилику за исправку и пут апелације. Излазак у хитној ситуацији треба да захтева јавни извештај о инцидентима и аутоматски преглед или залазак сунца.

\textbf{Забрањена листа.} Према овом шаблону, регистар не би допустио:

\begin{enumerate}[leftmargin=*]
\item инструменти који директно финансирају или стимулишу еколошко уништење изван уговољених граница, насиље или оружавање, присилно извлачење или системска злоупотреба; или
\item Клас ваучера који нема јасне одредбе о презентацији и испуњавању, одговорност и начине за исправљење.
\end{enumerate}
Список забрањених би био варијантиран, јавно ревидирани и мењао би се само кроз усвојен праг критичке акције и временски блок, приказан као Q3 + T3 у Додатку Д.

\subsection*{11.1 Разменни курс и ограничено управљање}
\addcontentsline{toc}{subsection}{11.1 Разменни курс и ограничено управљање}

\textbf{Промене у времену.} Разградња по овом шаблону би променила методе размене и одвојено имплементирала граничне параметре тек након јавне временске блокирања. Излазак у хитној ситуацији би користио одвојен процес одобрења и укључивао аутоматски запад сунца или преглед.

\textbf{Прагови одобрења.} У шаблону су предложени виши прагови одобрења за промене базе индекса вредности и глобалне промене граничног нивоа, усредни прагови за промене одређених трећих лица у групи и стандардни прагови за рутинске промене плате.

\textbf{Публиковани фед.} У учествујућим распоређивању би за сваку пулу објављивали променљиве индекса на ланцу, изворе или медијанте оракула, каденцију ажурирања, граничне прозоре и капане и режиме неуспеха или сигурне константе.

\textbf{Критерији за спешну паузу.} У учествујућим распоређивању би били предузведени услови као што су прекид оракула, висока ограничена употреба у комбинацији са неуспехама у испуњавању или неизменним неуспехом, заједно са проверима резюме и захтевима за преглед након инцидента.

\subsection*{Пример јавне индексне подаце за један Тржница и ваучер}
\addcontentsline{toc}{subsection}{Пример јавне индексне подаце за један Тржница и ваучер}

\begin{itemize}[leftmargin=*]
\item \textbf{Симбол:} на пример, \texttt{Maize\_50kg@IssuerY}.
\item \textbf{Референтна јединица:} Индекс јединица (IUX).
\item \textbf{Публикована вредност:} 30.000 IUX.
\item \textbf{Извор:} медијана идентификованих извора, као што су локална истраживања тржишта, обавештење министарства и база за распоређивање.
\item \textbf{Актуализовани каденција:} свакодневно у 18:00 EAT, са 24-часовим временским блокирањем.
\item \textbf{Режим неуспеха:} замрзнути на последњу валидну вредност, применити политику откривене границе и зауставити после 72 сати прекида.
\item \textbf{Рационално:} објављене ноте и запис о променама из претходног ажурирања.
\item \textbf{Подписници:} откривене мултисиг адресе и праг одобрења.
\end{itemize}
\subsection*{11.2 Предложена књига о осигурању}
\addcontentsline{toc}{subsection}{11.2 Предложена књига о осигурању}

\textbf{Само опционални дизајн.} Овај рунбук се односи само на распоређивање које је изрично усвојило и финансирало осигурални фонд и објавио покривене догађаје, прихватљиве захтевце, одговорно лице, средства, границе, искључења, захтеве за доказом, процес и услове управљања. Ни тренутни CLC App ни GEF не пружају покривеност само зато што се овај дизајн појављује у Белој књизи.

\textbf{Могуће изазове.} Прихваћена политика могла би да покрије неуповршавање дефинисаног емитента, недостатак резерве групе или губитак моста или гаранције. технички инцидент не квалификује се аутоматски; применљиве политике би контролисале.

\textbf{Проценка.} Одговорни орган би примирио квитанте за трансакције, залихе залиха, облигације гаранта, записи о подаци изкупа, одговоре емитента и друге потребне доказе, а затим би објавио извештај о инцидентима који је у складу са приватношћу и законом.

\textbf{Илустративни губитак водопада.} Уколико сваки слој постоји и законски се примењује, полица би могла користити: (1) одговорне облигације емитента или гарантске улоге $\to$ (2) резерве на нивоу пула $\to$ (3) предложен мрежни осигурачки фонд $\to$ (4) привремени смањење на захтев о опционалном покривању, само ако то изрично дозвољавају постојећи услови и примењиво законодавство $\to$ (5) законско повлачење за доказану превару или злоупотребу.

Корекција покривености не смањује основне обавезе емитента о ваучеру или не мења баланс на ланцу, осим ако то изричито дозвољавају важећи већ постојећи услови и примењиво законодавство и ако се добије било која потребна согласност власника.

\textbf{Ограничења и искључења.} Објављено покривање би дефинисало ограничења, прихватљиве презентације, доказе, прозоре захтева, искључене маршруте или догађаје, географске ограничења и третирање исцрпљених резерва. Изплата би могла бити нула након што се достигну границе које се примењују.

\textbf{Илустративни распоред за опоравак.} Ако је усвојена и објављена:

\begin{enumerate}[leftmargin=*]
\item позиве би прво износиле из одговорног емитента или гаранта облигација, затим из одговарајућих резерва резерви, а затим из предложеног мрежног осигураног фонда;
\item било које смањење на захтев о опционалном покривању би било ограничено на оно што су већ постојећи услови покривања и примењиво законодавство дозвољавали, до објављеног ограничавања инцидента;
\item план за опоравак би могао применити наведен део опорављене вредности за наведен период, након чега би се сваки преосталан покривен недостатак претворио у регистровани губитак са јавним постмортем; и
\item Свака одлука би пружила квитанцу са инцидентом ID, повреденим захтевима и ваучерима, одлуком, планом за опоравак и прозором апелације.
\end{enumerate}
\subsection*{11.3 Рамка гаранције}
\addcontentsline{toc}{subsection}{11.3 Рамка гаранције}

У овом одељку се разликује одговорност емитента, опционалне заштите фондова и гаранције трећих страна. Пулеви могу да се такмиче за курацију, услове и изричито понуђене заштите без подразумевања да CLC App, CPP, GEF, или било која ширија мрежа аутоматски гарантују ваучер.

\subsection*{Одговорност исходног емитента}
\addcontentsline{toc}{subsection}{Одговорност исходног емитента}

\begin{itemize}[leftmargin=*]
\item Сваки ваучер је пре свега одговорност издаваоца. Издајца се обавезује да обезбеђује пријављену робу, услугу или законски новчани еквивалент на својим објављеним условима.
\item Издавачи би објављивали ко може да представи ваучер, шта значи испуњење, где и када је доступан, који докази су потребни и који се лекови примењују.
\item Ако емитент не испуни захтев, емитент је главни одговорна страна. Заштите резервоара или мрежа се примењују само када се одвојено усвоје, финансира и открије.
\end{itemize}
\subsection*{Опционална заштитна опрема}
\addcontentsline{toc}{subsection}{Опционална заштитна опрема}

Управник Тржнице може изабрати да додаје ограничено дефинисану заштиту допуњеним ваучерима. Она није аутоматска и требало би да идентификује одговорну страну, финансирање, прихватљиве догађаје, ограничења, прозоре, доказе, искључења и средства у метадатовима пула и примењивим условима.

Илустративни типови заштите укључују:

\begin{enumerate}[leftmargin=*]
\item \textbf{Покривање резервних средстава:} Након потврђеног неиспостављања емитента, одговорни субјекти групе плаћају дефинисан износ у одређеном резервном активу, уз обзир на своју објављену границу и доступне финансиране резерве.
\item \textbf{Процена за повраћање:} Након квалификованог догађаја, група нуди временски ограничени пут за суап у претходно одобрен актив или у други одобрен актив, под обзиром на ограничења и инвентар. Ово је заштита ликвидности која зависи од инвентара, а не обећање да је свака суап реверзивна.
\item \textbf{Алтернативна испуњавања:} одговорна страна организује одобреног заменног пружаоца у пределу објављеног ограничења количине или вредности.
\item \textbf{Заштита размене:} за изабране класе ваучера, група нуди само прилагођавање покривености или средство за суап-бацк који је наведено у својим већ постојећим условима. То не смањује основну обавезу емитента о ваучеру.
\end{enumerate}
\subsection*{Могући извори финансирања}
\addcontentsline{toc}{subsection}{Могући извори финансирања}

\begin{itemize}[leftmargin=*]
\item \textbf{Облигација емитента:} обезпеке које је емитент поставио или задржао у откривеном резерву и које су доступне након потврђеног покривеног догађаја.
\item \textbf{Резерва Тржница:} имовине које контролише одговоран субјект "Пол" и које су додељене заштитима које оглашава.
\item \textbf{Облигације гаранта треће стране:} обезпеке које је увео идентификован спољни гарант за пријављене емитенте, класе ваучера или догађаје.
\end{itemize}
У учеству гаранта би се следили објављени критеријуми прихватљивости, величина облигација, границе концентрације, овлашћење доношења одлука и правила за законско спровођење.

\subsection*{Процес захтева}
\addcontentsline{toc}{subsection}{Процес захтева}

У уведеној политици би се дефинисале проводнице које се могу ревидирати, као што су рок испуњавања који је пропуштен након ваљиног податка изкупа, потврђена неплачност емитента, покривен мост или неуспех гаранције, или формално декларисано стање инцидента. Такође би дефинисала:

\begin{itemize}[leftmargin=*]
\item Како учесник отвара захтев и пружа потребне доказе о представљању и испуњавању;
\item који проверава услове ваучера, одговоре емитента и техничке записи;
\item одлуку и прозоре о жалби; и
\item дозвољен пут плаћања, средства, ограничења и квитанце.
\end{itemize}
Укупни приходи од емитента, арбитража или законског спровођења би поновили важеће облигације или резерве у складу са објављеном политиком пре него што би се користили за предложен приступ CLC мрежном пулу суапа.

\subsection*{Потребно откривање}
\addcontentsline{toc}{subsection}{Потребно откривање}

За сваку покривљену категорију пула и ваучера, одговорна страна би објавила:

\begin{itemize}[leftmargin=*]
\item да ли је гарант одсутан, изборни или захтеван;
\item величина облигација или резерва и ограничења концентрације;
\item типове заштите, средства, капане, прозоре и искључења;
\item Срок за представљање, испуњавање, захтев и апелацију; и
\item Једноставна изјава о томе ко гарантује шта и шта није гарантовано.
\end{itemize}
\textbf{Принцип курације.} Управници Тржница и одговорне правне или управљачке структуре одговорне су за заштиту коју рекламирају. CPP- компатибилно распоређивање може да обезбеди стандарде, регистре или опционалне заједничке политике, али ни CLC ни GEF аутоматски не гарантују купоне или пулове.

\subsection*{11.4 Заштит против ухвативања}
\addcontentsline{toc}{subsection}{11.4 Заштит против ухвативања}

Према овом шаблону, следеће би биле критичне акције које захтевају усвојен највиши ниво одобрења и дуго време:

\begin{enumerate}[leftmargin=*]
\item промена предлажена плата, укључујући покривеност и приоритете главних операција;
\item промена корена канонског регистра;
\item промена опсега покривености, границе захтева или овлашћења за доношење одлука;
\item проширење моћи за паузу у хитној ситуацији; или
\item ослањавање обавеза о форкибилности, транспарентности или суверенитету поља које су наведене у овом документу.
\end{enumerate}
\subsection*{11.5 Процедура излаза и излаза}
\addcontentsline{toc}{subsection}{11.5 Процедура излаза и излаза}

Ако је управљање ухваћено или вредности се значајно одвијају, заједнице, Управници Тржница и оператори би могли покушати да изађу од тога крскајући слој управљања мреже. Континуитет основних пула и купонова зависи од распоређених уговора, кључева, интерфејса, инфраструктуре, услуга трећих страна и применљивих обавеза.

Процес напуштања би могао:

\begin{enumerate}[leftmargin=*]
\item \textbf{Публикујте снимку:} износи изабране регистре, ваучере, вредности, границе и политику плате, а затим објављује потписан хаш снимка.
\item \textbf{Реинсталирајте услуге управљања:} унедрити нове корене регистъра, услуге маршрута и усвојене модуле таксе или покривености у новој структури одговорности.
\item \textbf{Презапис:} дозвољавати Управници Тржница да се пријави регистровањем своје адресе групе под новим кореном, без захтева од власника да мигрирају ваучере који иначе функционишу.
\item \textbf{Репоновани клијенти} додати нови корен као изборни мрежни профил у СДК-ма и интерфејсима, са било којим дефолтним променама направљеним кроз раскључен процес управљања.
\item \textbf{Управљајте мостовим временом:} одржавање компатибилних путева где је сигурно и одбијање путева који крше правила новог профила.
\end{enumerate}
Циљ пројекта је да напуштање канонског регистра не онемогућава другачије функционалне локалне базе. Стварна континуитет остаје зависна од распореда; Федерација је слој откривања и координације.


\section*{12. Предложена ликвидација програма ликвидности (незавршавајући обхват)}
\addcontentsline{toc}{section}{12. Предложена ликвидација програма ликвидности (незавршавајући обхват)}

Ово је илустративна контролна листа за будући, одвојено документовани програм. То није понуда, тренутна функција апликације или обећање ликвидности, приступа, осигурања, повратака, вредности или исхода.

\begin{itemize}[leftmargin=*]
\item \textbf{Прихватљиви доприноси:} Стабилни монети, резервни токен или одобрени купонови заједнице, као што је наведено у програму.
\item \textbf{Постављање:} одређено Тржнице обавеза или предложен CLC мрежни Тржница у оквиру одвојеног мандата.
\item \textbf{Закључавање и излазак:} Програмски специфични услови, као што су периоди од 30 до 90 дана и откривене врата под стресом.
\item \textbf{Кредит за таксе које су постављене у политику:} опционални механизам "ex-post" под објављеним капама и прозорима, а не обећано повратак.
\item \textbf{Дељење ризика:} било које смањење на захтев о опционалном покривању мора бити изрично одобрена већ постојећим условима и примењивим законом. Она сама не мења обавезу емитента о ваучеру или баланс на ланцу.
\item \textbf{Извештај:} периодичне дашбодове и потврде који покривају здравље Тржница, инвентар, границе, таксе и све финансиране резерве покривености.
\item \textbf{Споразуми:} ограничења коришћења прихода, контроле оракула, граничне нивое, правила конфликата и средства за опоравак.
\item \textbf{Могућа прихватљивост предложених токова:} доприносиоци који поседују одређене пулове могу бити квалификовани за поседу предложених CLC грантова за управљање токонима на основу мереног приступа пулу-свају и потврђеног испуњавања емитента, уз обзиром на усвојена услова, правила против игара и ограничења политике.
\end{itemize}


\section*{13. Глосар на једноставном језику}
\addcontentsline{toc}{section}{13. Глосар на једноставном језику}

\begin{itemize}[leftmargin=*]
\item \textbf{Cosmo-Local Credit (CLC):} Семейство производа, укључујући CLC App, документацију и шири рад о заједничкој ангажовању.
\item \textbf{CLC App:} Прогресивна веб апликација коју управља GEF на \texttt{cosmolocal.credit}.
\item \textbf{Протокол о заједничкој обавези (CPP):} Повторно коришћени модел курације, процене, ограничења, размене и одговорног управљања.
\item \textbf{Protocol v1.1.0:} Проверене јавно објављивање паметних уговора.
\item \textbf{Тржница обавеза:} Уред који се регулише који прихвата подржане имовине и објављује правила размене.
\item \textbf{\texttt{SwapPool}:} Тренутни симболски трезор и директни договор за размену.
\item \textbf{Знак:} Баланс на ланцу или дигитални актив. Механика токова сама не ствара обећање које се може изкупити.
\item \textbf{Купон:} Токен или запис који емитент представља као обавеза која се може изкупити на објављеним условима.
\item \textbf{Подношење:} Каталог записи за стоке или услуге повезане са ваучером.
\item \textbf{издавалац:} Страна одговорна за објављивање и испуњавање обавезе о ваучеру.
\item \textbf{Управник Тржнице:} Одговорна структура која објављује и управља правилима Пола.
\item \textbf{Обменни курс или котирање:} Параметр трансакције који је произведен конфигурисаним котирачем; није доказ готовности или вредности откупа.
\item \textbf{Покупни токен-баланс капа:} Тренутни \texttt{Limiter} максимум за токен у власништву пула. Апликација може да га означи \textbf{ Граница кредита.''}
\item \textbf{Измена резервоара:} Директна размена средства преко једног \texttt{SwapPool}.
\item \textbf{Уплата за суап:} Трансфери у ланцу и рачуноводство плате који заврше суап пула.
\item \textbf{Презентација о откупу:} Враћање или другачије представљање ваучерских јединица емитенту.
\item \textbf{Извршавање:} Емитент пружа обећано добро, услугу, корист или другу резултатност.
\item \textbf{Увођење:} Допис који спречава повторну презентацију испуњених јединица.
\item \textbf{Предложен рутер за извршење:} Будући софтвер који би одобрио и извршавао маршруте. Тек тренутни \texttt{SwapRouter} цитати.
\item \textbf{Предложена мрежа CLC:} Будући уговор о клирингу и координацији буџета на нивоу мреже.
\item \textbf{Предложена CLC токен управљања:} Будући дизајн управљања, а не тренутни Апп или Актив Protocol v1.1.0.
\item \textbf{Рек сетке:} Предложен удео таксама за покупни фонд који се преноси на будући буџет за мрежу.
\item \textbf{Гаранција или осигурање:} Заштита која се примењује само када је одређена страна изрично преузима, финансира и објављује.
\end{itemize}

\section*{14. Предложени КПИ и показатељи здравља}
\addcontentsline{toc}{section}{14. Предложени КПИ и показатељи здравља}

Разградња треба да објављује KPI само када може дефинисати догађај, извор, кохорту или период, јединицу, методу процене и временски печат, искључења, политику исправљања, доказе изван ланца и ограничења квалитета података.

Предложене мере укључују:

\begin{itemize}[leftmargin=*]
\item важећи подаци о откупу;
\item стопа испуњавања на основу кохорте;
\item потпуност испуштања;
\item латенција од презентације до испуњавања;
\item трајање задржавања од прикупљања до презентације;
\item изложено допустимо обавеза у складу са наведеном методологијом;
\item завршен обем суапа резервоара;
\item мерена инвентарска група;
\item Употреба токена-баланса капитала;
\item скорост проласка котирања, одржана одвојено од скорости извршења маршрута;
\item допустима резерва подељена по експлицитно покривеној експозицији;
\item позиве гаранта и повраћања у оквиру идентификоване политике;
\item предложени мрежни удео и накнаде за услуге који су стварно примљени, без двоструког рачунања бруто накнада Тржница;
\item време за откривање управљања, одлуку, паузу, поправку и затварање;
\item жалбе, спорове, исправке и средства за опоравак; и
\item Одвојени су социјални или еколошки исходи.
\end{itemize}
Подаци о трансакцијама на ланцу само по себи не доказују идентитет, перформансе емитента, задовољство, причинност, здравље заједнице или социјални утицај. Ове тврдње захтевају сопствену методологију и доказе.

Видите ли \href{https://docs.cosmolocal.credit/sr/white-paper/appendix-c-kpi-definitions}{Додатко С} за предложену спецификацију мерења.


\section*{15. Дорачна карта (индикативна)}
\addcontentsline{toc}{section}{15. Дорачна карта (индикативна)}

\subsection*{15.1 Тренутни темељ}
\addcontentsline{toc}{subsection}{15.1 Тренутни темељ}

GEF управља CLC App на \texttt{cosmolocal.credit} на ланцу гнозе. Апликација пружа поддржаван приступ рачуну и новчанику, јавни каталог тржишта и интерфесе за токове, ваучере, понуде, Тржнице обавеза, трансфери и директне суапе Поул.

Protocol v1.1.0 обезбеђује основу садашњег уговора: \texttt{GiftableToken}, директна извршавање \texttt{SwapPool}, изборни регистри, модули за процјену, ограничења баланса токена пула, таксе пула, додатна такса протокола и само цитат \texttt{SwapRouter}. Доступност још увек зависи од интерфејса, распоређених уговора, инвентаризације, конфигурације, државе мреже, допустимости, пружаоца, надлежности и објављених услова емитента или пула.

\subsection*{15.2 Предложени мелистонови}
\addcontentsline{toc}{subsection}{15.2 Предложени мелистонови}

Ови названи етапи су напутства за дизајн, а не бројеви објављивања, датуми испоруке или обавезе које ће функција покренути.

\begin{itemize}[leftmargin=*]
\item \textbf{Фондација за управљање:} предложени токен управљања CLC; Предложена мрежа CLC; прилагођавачи накнада; кворум, временски ограничење и темеље управљања.
\item \textbf{Направљање и посматрање:} Рутер SDK и регистарске АПИ; здравствени тачбодови; предложени оквир полице осигурања; намере на ребалансирање; прототип за сечење партова.
\item \textbf{Инструменти за ризик између домена:} HTLC или маршрутизацију по покровитељству; одвојени модули гаранта; Редактиви, рачуна и нивоирани гранични подаци.
\item \textbf{Регулиран приступ:} услуге плаћања трећих страна специфичне за распоређивање; личне микробачеве; откривање услуге у складу са стандардом; ревизије класа ваучера од стране трећих лица.
\end{itemize}
Сваки платежни сервис пружају одвојено идентификоване треће стране у складу са применим јурисдикцијом, прихватљивошћу, таксама, границама и условима. Уговор CLC App и Protocol v1.1.0 сами не користе фиатске реке.

Извршење мулти-хопа, заједничко осигурање, предложени CLC управљачки токен и предложен CLC мрежни пул, сетчано сетирање, лични микропулови и регулисани платежни сервиси остају предложени или зависни од распоређивања док их имплементација и њени правила не идентификују као активне.


\section*{16. Предложена шаблона процене}
\addcontentsline{toc}{section}{16. Предложена шаблона процене}

Регистар, пул или будући програм ликвидности могу прилагодити овај шаблон. То није тренутни универзални стандард за листирање или гаранција.

\begin{enumerate}[leftmargin=*]
\item \textbf{Набљуђене чињенице:} Идентичност и историја емитента, услови ваучера, понуде, доказ капацитета, конфигурација пула, контролера, инвентар, границе, резерве, гаранте, инциденти и нерешене обавезе.
\item \textbf{Циљ и погођене стране:} намењену корист, ризике, конфликте, еколошке и друштвене ефекте и ко носи губитак или одговорност.
\item \textbf{Оценка:} Цена понуде, показана вредност ваучера, метода разменних стопа, референција оракла, јединице, временске штампе и разлози за било какву дивергенцију.
\item \textbf{Границе:} Прихватљивост, забрањене употребе, концентрација, границе равнотеже токена, паузе, географске или правне ограничења и покретачи за преиспитивање.
\item \textbf{Заштита:} Идентификована обавезна страна, покривено догађај, финансирање, ограничење, искључења, трајање, доказ, процес захтева и средства за опоравак.
\item \textbf{Одлука:} одобре, ревидира, одбацује, паузује или ескалује за људску ревизију; Идентификујте доносец одлука, разлоге, услове, датум прегледа и процес апелације или исправљања.
\end{enumerate}
Излаз не би требало да описује курацију, границе, резерве или верификацију као одобрење, капацитет емитента, осигурање или гарантовано испуњење, осим ако одговорна страна изрично није утврдила ову заштиту.


\section*{17. Правни и обавештења о поштовању}
\addcontentsline{toc}{section}{17. Правни и обавештења о поштовању}

CPP може координисати токове, ваучере, пулове и суапе, чији законски третман зависи од њиховог дизајна, маркетинга, употребе, одговорних страна и надлежности. Ништа у овој Белој књизи није правно, пореско, инвестиционо, кредитно или финансијски савет.

Употребу CLC App уређују \href{https://docs.cosmolocal.credit/sr/governance/terms}{Услови коришћења}. GEF управља апликацијом и пратећом инфраструктуром. Осим ако GEF изричито не преузме другу улогу, није издавалац, Управник Тржнице, чувар имовине, зајмодавац, зајмопримац, посредник, гарант, страна која искоришћава ваучер, саветник, осигуравач нити страна у обавезама између учесника.

Службе плаћања које зависе од распореда морају идентификовати свог одговорног пружаоца, юрисдикцију, прихватљивост, таксе, границе, модел чувања и услове. Присуство повезаника не значи да GEF или Поул управља основном регулисаном послу.

\subsection*{17.1 Појаве о ваучерима и понудама}
\addcontentsline{toc}{subsection}{17.1 Појаве о ваучерима и понудама}

Емитент треба да објављује и одржава актуелност:

\begin{itemize}[leftmargin=*]
\item одговорни идентитет и контактне информације;
\item прикључена понуда, јединица, снабдевање, наведена вредност и капацитет;
\item локације и процедуре презентације;
\item време испуњавања и доказ;
\item правила за истечење, таксе, порезе, ограничења и промену материјала;
\item жалбе, замење, повраћај или друге средства за опоравак; и
\item метод испуштања који спречава повторну употребу након испуњавања.
\end{itemize}
Токенски уговор не утврђује ове услове нити доказује извршење.

\subsection*{17.2 Публикације резервоара}
\addcontentsline{toc}{subsection}{17.2 Публикације резервоара}

У Управник Тржнице треба да се објављују:

\begin{itemize}[leftmargin=*]
\item сврха базна и одговорни доносељи одлука;
\item власник \texttt{SwapPool}, прокси администратор, контролер зависности и примаоци плате;
\item допуњени средства и правила суспензије или уклањања;
\item методе процене, изворе котирања, таксе, границе баланса токена, инвентар и овлашћења за повлачење власника;
\item права доприноса и излаза;
\item све резерве, гаранције, гаранте, полисе осигурања и правила о расподелу губитака; и
\item управљање, унапређење, хитне ситуације, жалбе, миграције и процеси прекида.
\end{itemize}
Листирање није одобрење, процене, осигурање или гаранција GEF.

\subsection*{17.3 Акције и обавезе}
\addcontentsline{toc}{subsection}{17.3 Акције и обавезе}

Обични човек. \textbf{Измена базе} Биржа подржавају имовине у складу са приказаном котирањем и трансакционим границама. Управљање имовином не чини да је кредит, исплата, искупљање емитента или реализација у стварном свету.

\textbf{Презентација о откупу } враћа или приказује јединице ваучера емитенту. \textbf{ Испуњење } је обећани резултат емитента. \textbf{ Раздужење } записи испуњене јединице тако да се не могу поново користити. Апликација је \textbf{„Искористи``} акција тренутно припрема трансфер власнику токена; само тај пренос не доказује испуњавање.

Апликација је \textbf{„Повуци ваучер``} Акција је реверзивна искључивање каталога. Она не спаљава балансе, не отказује захтеве или не испуњава обавезе.

Будући кредит или други кредитни производ би захтевао одвојене додатне и трансактивне услове пре употребе. Нису обични послати, доприносити, депозити у пулу, суапи у пулу, презентације, испуњавања или разплате који стварају кредит само због његовог услова или типа средства.

\subsection*{17.4 Заштита, докази и локални закон}
\addcontentsline{toc}{subsection}{17.4 Заштита, докази и локални закон}

Лимит, резерва, упис у регистар, листирање апликација или блокчејн трансакција нису аутоматски гаранција или осигурање. Свака заштита мора да идентификује одговорну особу, покривено догађај, финансирање, ограничење, искључења, трајање, доказ, процес захтева и важеће услове.

Доказања из јавне ланце могу показати адресе, имовине, суме, часове и догађаје у договору. Само по себи не доказује идентитет, способност издаваоца, испуњавање, задовољство, разматрање управљања, правно ослобођење или социјални утицај.

Функције могу бити ограничене или недоступне због закона, санкција, прихватљивости, државе уговора, инвентаризације, граница, безбедности, јурисдикције или услуга треће стране. Издајци, Управници Тржница, пружаоци и учесници остају одговорни за одређивање и поштовање примењивог закона.


\section*{18. Закључак}
\addcontentsline{toc}{section}{18. Закључак}

Cosmo-Local Credit повезује актуелну апликацију и Protocol v1.1.0 темељ са ширијим предложеним дизајном за независно управљање Тржнице обавеза. Тренутни директни суапи пула, компоненте котирања и лимита и јавни записи трансакција могу подржати одговорну размену, али не доказују испуњавање емитента, не стварају права доприносача или гарантују вредност, ликвидност, искоришћење, осигурање, правни статус или заштиту од губитака.

Предложено рутирање извршења, нетринг, управљање имовином, клерирање мреже, програме ликвидности и заједничке заштите описане у овом документу захтевају одвојене имплементације, финансирање, управљање и објављене услове. Конструкција има за циљ да се прошири интерперабилност док се одржава перформанса емитента, управљање пулом и одговорност у стварном свету са идентификованим странама.


\section*{А. А. Предложен модел мерења и плате}
\addcontentsline{toc}{section}{А. А. Предложен модел мерења и плате}

У овом додатку се дефинише предложен оквир мерења. Protocol v1.1.0 не записва испуњавање емитента, испуњење у стварном свету или свако поле података које се захтева испод.

\subsection*{Дефиниције догађаја и залиха}
\addcontentsline{toc}{subsection}{Дефиниције догађаја и залиха}

За класу ваучера \emph{j}, кохорту или период \emph{t}, и откривену методу процене \emph{m}:

\begin{itemize}[leftmargin=*]
\item \texttt{O\_\{j,t,m\}}: вредност непоплаћених прихватљивих обавеза на граници мерења.
\item \texttt{X\_\{j,t,m\}}: вредност завршених суапања групе током периода.
\item \texttt{P\_\{j,t\}}: јединице које се валидно представљају емитенту за искоришћење.
\item \texttt{F\_\{j,t\}}: представљене јединице са одвојено доказаном испуњавањем емитента.
\item \texttt{G\_\{j,t\}}: испуњене јединице са записом испуштања које спречавају повторну употребу.
\end{itemize}
\texttt{O} се не може закључити само из понуде токена. У политици мерења треба да се идентификује одговорни емитент и да се, по прикладностима, искључе инвентар у власништву емитента, спаљене јединице, истеле јединице, ослобођене јединице, тестове балансе, недоступне балансе и токове чији услови не стварају непостојану обавезу треће стране.

Свака мера која се процењује мора објављивати јединицу, извор, методу процене, временски печат и третман разних обменних курсева полу. Трансфер на ланцу може подржати \texttt{X} или доказ о презентацији; само по себи не успоставља \texttt{F} или \texttt{G}.

\subsection*{Мерке испуњавања засноване на кохорти}
\addcontentsline{toc}{subsection}{Мерке испуњавања засноване на кохорти}

За кохорту важећих претстава о изкупу:

\texttt{FulfillmentRate = FulfilledPresentments / ValidPresentments}

\texttt{DischargeCompleteness = DischargedFulfillments / FulfilledPresentments}

У сваком бројевнику и знаменачу користите исти затворен или зрели кохорт. Извештај одбачен, повлачен, истекао, оспорен, делимично испуњен, исправљен и још увек отворени подаци одвојени.

\texttt{FulfillmentLatency = median(t\_fulfilled - t\_presented)}

\texttt{HoldingDuration = median(t\_presented - t\_acquired)}

Извршавање латенције мере услугу емитента након презентације. Трајање задржавања је посебна мера и не сме да се означи латенција откупа.

\subsection*{Различне мере брзине}
\addcontentsline{toc}{subsection}{Различне мере брзине}

Предложена брзина испуњавања обавеза може се израчунати само ако \texttt{O} и испуњена вредност користе исте методе процене:

\[
V_{\mathrm{commit},t} = \frac{\mathrm{FulfilledValue}_t}{\mathrm{AverageOutstandingEligibleCommitmentValue}_t}
\]

Измер за суап-активност пула је одвојен:

\texttt{V\_swap,t = PoolSwapValue\_t / AverageMeasuredPoolInventoryValue\_t}

Ни једна од ових вредности не доказује социјални утицај, капацитет емитента, профитабилност или конвертабилност новца.

\subsection*{Предложени приход од мреже}
\addcontentsline{toc}{subsection}{Предложени приход од мреже}

Нека:

\begin{itemize}[leftmargin=*]
\item \texttt{PF\_t} су бруто пошлине од групе генерисане током периода;
\item \texttt{NR\_t} је предложени мрежни грејк који је заправо добио као откривен део тих плате по групи;
\item \texttt{RF\_t} су одвојене предложене трошкове за маршрутизацију или услуге које су заправо примљене; и
\item \texttt{$\chi$\_t} је измењен део прихода који је прихватљив и конвертибилан за декларисану употребу у новчаној номинацији након трошкова и ограничења политике.
\end{itemize}
Из предложеног мрежног буџета:

\texttt{ProposedNetworkRevenue\_t = NR\_t + RF\_t}

\texttt{CashUsableNetworkRevenue\_t = $\chi$\_t $\times$ (NR\_t + RF\_t)}

Не додајте \texttt{PF\_t} на \texttt{NR\_t}: раке је превод од брутог трошкова Пул и иначе би се рачунао два пута. тренутни Protocol v1.1.0 уместо тога подржава додатну таксу протокола; Његови квитанси морају бити пријављени одвојено од овог предложеног модела раке.


\section*{Б. Иллюстративна будућа плата водопад}
\addcontentsline{toc}{section}{Б. Иллюстративна будућа плата водопад}

Ово је предложен дизајн за будуће распоређивање. Она се примењује само ако се имплементира, финансира и усвоји путем објављеног управљања и условима учесника. У њему се не описује тренутна такса Protocol v1.1.0, право доприносника, резерва, гаранција или осигурање.

Нека \texttt{F\_in} буде приход који је заправо добио предложен буџет за мрежу током епохе: предложену мрежу, одвојене трошкове за маршрутизацију или услуге и било који други изрично одређени приход. Укупне плате за Тржница који остају са Тржницом нису укључене.

Активи прихода морају се класификовати као:

\begin{itemize}[leftmargin=*]
\item \texttt{E\_cash}: активи који су допустиви за декларисану употребу у готовности у складу са усвојеном политиком; или
\item \texttt{E\_kind}: ваучер у природи или други активи који се не третирају као конвертибилни у готовини.
\end{itemize}
У свакој политици конверзије би требало да се региструју средства и места, одговорне власти, изворице цијена, границе проклањањања, извештавање и важеће правне контроле.

Предложен водопад може применити приход који је добио у овом редоследу:

\begin{enumerate}[leftmargin=*]
\item \textbf{Циљ покривене резерве:} финансирање усвојена резервна или осигурана мета заснована само на дефинисаним покривеним експозицијама, прихватљивим имовинама, искључењима и правилима захтева.
\item \textbf{Основне операције:} финансирање откривеног, ограниченог оперативног буџета.
\item \textbf{Мандати за ликвидност:} Фонд је одобрио, одвојено управљао пуле или маршрутизационим програмима.
\item \textbf{Остатак буџета:} доделити остатак износа између операција, програма ликвидности и додатног буфера под објављеним ограничењима.
\end{enumerate}
Циљ резерве сам по себи није гаранција. У сваком осигурању или гаранцији треба да се идентификује обавезанца, покривеног догађаја, финансирања, ограничења, искључења, трајања, доказа, процеса захтева и доделе губитака.

\textbf{Гвардија:} Подељке за водопада су намењене за усвојене покривености, операције и услуге ликвидности. Они се не смеју обликовати или извршити као операције цене подршке.

Према предложеном моделу, било који предложен CLC токен за управљање који се стече путем посебно омогућеног програма спољне ликвидности би се пензионисао или стављао у откривену лавицу без гласања. Овај механизам није распоређен.


\section*{Ц. Предложене дефиниције KPI}
\addcontentsline{toc}{section}{Ц. Предложене дефиниције KPI}

Ови КПИ чине предложену спецификацију мерења, а не изјаву да тренутна апликација или протокол записују сва потребна догађаја.

Сваки објављен KPI мора да садржи:

\begin{itemize}[leftmargin=*]
\item одговорни издавач и извор података;
\item дефиниција догађаја или државе;
\item кохортни период или период мерења;
\item Јединица и метода процене;
\item временски штамп за процене;
\item правила укључивања и искључивања;
\item обраду делимичних, спорних, истекалих, неприступних и исправљених података;
\item захтеван доказ или потврда изван ланца;
\item историју ревизије и ограничења квалитета података; и
\item да ли је резултат на ланцу, пријављен, потврђен, независно потврђен или процењен.
\end{itemize}
\begin{longtable}{P{0.31\linewidth} P{0.31\linewidth} P{0.31\linewidth}}
\toprule
KPI & Предложена дефиниција & Потребни докази и искључења \\
\midrule
\endhead
\textbf{Валидна презентација} & Број или јединице које су прихваћене у процес искупљања емитента током периода & Идентификатор презентације, емитент, овлашћење власника, износ, време, статус; искључити дуплите и неважеће захтеве \\
\textbf{Стопа испуњавања} & Завршени валидни подаци подељени на валидни подаци за исте зреле кохорте & Одвојени докази о резултатима емитента; извештавање о отвореним, одбаченим, спорним, делимичним и исправеним случајевима \\
\textbf{Површеност испуштања} & Завршени подаци са записом ослобођења подељени по испуњеним подацима & Пожар, укидање, запис о онемогућавању или други доказ о непоновном употреби повезан са испуњавањем \\
\textbf{Латенција испуњавања} & Медијан и 90-ти процентил времена од презентације до испуњавања & Не замењују време задржавања емисије до презентације или прикупљања до презентације \\
\textbf{Трајање одржавања} & Медијан и дистрибуција времена од прикупљања до подавања & Идентификујте догађај прикупљања и искључите непознате времена прикупљања \\
\textbf{Остале прихватљиве обавезе} & Прихватљиве обавезе треће стране које остају након дефинисаних искључења & Идентичност емитента и услови; искључити важећи инвентар емитента, излагање, спаљење, раздужење, тестирање и токене који нису обавезни \\
\textbf{Област обмена Тржницом} & Стоимост завршених директних суапа група у складу са једном откривеном методом процене & Процес у вези са уређивањем на ланцу, јединици токена, извор стопе, време; не класификују као испуњење емитента \\
\textbf{Инвентар Тржница} & Измерени подржани средства у власништву пула у одређеном времену & Контрактни баланси, резерви за таксе, неприступни средства, метода процене и овлашћења за повлачење власника \\
\textbf{Одличност резерви} & Добиљни, доступни резервни средства подељени по изрично покривеној експозицији & Политике покривености, допустимост актива, чување/контрола, пасиве, искључења и процене; не укупна испорука токова по дефолту \\
\textbf{Лимит коришћења} & Измерени баланс токена пула подељен на тренутно конфигуриран капа & Limiter адреса, токен, пул, временски печат, промене и периоди без ограничавајућег \\
\textbf{Цотира пролаз стопа} & Успешни одговори на цитати подељени валидним покушајима цитата & Резултат је само цитат; не извршавање маршрута \\
\textbf{Скорост извршења маршрута} & Завршени извршени мулти-хоп подељени валидним покушајима извршења & примењује се само на имплементисан систем извршења; Извештај по правилима по хопу и атомичности \\
\textbf{Опоравак гаранта} & Добијена прихватљива повраћања подељена по исплаћеном покривеном позивом & Идентификовани гарант, политика захтева, распоређивање, трошкови, спорови и списања \\
\textbf{Предложени приход од мреже} & Предложено је добијање мрежног ракета плюс добијање одвојених трошкова за маршрутизацију/услугу & Исклучује се бруто таксе на Тржници које задржавају Тржници и избегава се два пута рачунање ракета \\
\textbf{Временност управљања} & Време за откривање, одлуку, паузу, поправку и затварање инцидента & Опредељени часовници, одговорни органи, хитне овлашћења, апелације и нестали догађаји \\
\bottomrule
\end{longtable}

Позиве о друштвеним утицајима захтевају посебну методологију. Блокчејн активност сама по себи не утврђује идентитет, перформансе емитента, задовољство, здравље заједнице, причинност или утицај.


\section*{Д. Предложени параметри лансирања}
\addcontentsline{toc}{section}{Д. Предложени параметри лансирања}

Свака вредност испод је иллюстративна. То није тренутна дефолтна апликација, Protocol v1.1.0 подешавање, гаранција, понуда или обавеза испоруке. Будуће распоређивање би требало да усвоји, спроведе, прати и објави своје стварне параметре.

\begin{itemize}[leftmargin=*]
\item \textbf{Корумни нивоа за предложену CLC токену управљања:}
\begin{itemize}[leftmargin=*]
\item Рутина Q1: најмање 4\% кворум и више од 50\% одобрења.
\item Q2 осетљиво: најмање 10\% кворум и најмање 60\% одобрење.
\item Q3 критичан: најмање 20\% кворум и најмање 66.7\% одобрење.
\end{itemize}
\item \textbf{Предложено време:}
\begin{itemize}[leftmargin=*]
\item T1: 48 сати за Q1 акције.
\item T2: 7 дана за Q2 акције.
\item T3: 30 дана за Q3 акције.
\end{itemize}
\item \textbf{Каденца епоха:} 7 дана, укључујући и било које усвојене гласање, објављивање буџета и предложене прозоре sCLC.
\item \textbf{Изненадна пауза:} непосредно преко идентификованог органа за хитне ситуације, који истече након 72 сата, осим ако није ратификовано у складу са усвојеним процесима.
\item \textbf{Предложена мрежа:} 20\% такса за учествовање у групи по неизбежности, ограничена политиком.
\item \textbf{Илустративни опсег плате за Тржница:} 0\%--20\%, које објављује свака учесница.
\item \textbf{Предложено ограничење трошкове за маршрутизацију или услуге:} 20 БПП по маршруту.
\item \textbf{Предложена стопа мрежа:} \texttt{rake\_rate\_p = pool\_fee\_p $\times$ rake\_share\_p}.
\item \textbf{Прихватљивост прихода:} класификују примљене средства као готовне средства\texttt{E\_cash}или у природи\texttt{E\_kind}по објављеној методи.
\item \textbf{Контроле конверзије:} Дозволници средства и места, одговорни органи, извори цена, временски прозори, границе проклањања, ограничења и извештавање.
\item \textbf{Буџетска дозвола:} Публикују предложен буџет за приступ таксима само након што су усвојени покривени резерви и оперативни циљеви испуњени; буџет може бити нула.
\item \textbf{Линије ризика:} не излагају једну категорију средства ризику друге категорије без експлицитног, информисаног приступа у складу са применљивим законима.
\item \textbf{Предложени ограничења за ролинг и рачуна:} одбијају неразвијену транскласну активност и примењују усвојена ограничења.
\item \textbf{Опционално смањење покривености:} само ако то дозвољавају већ постојећи услови покривања, потребна согласност и примењиво законодавство; само по себи не смањује обавезу емитента о ваучеру или баланс у ланцу.
\item \textbf{Извънне контроле ликвидности, ако су посебно омогућене:} објављују опсег места, методу мерења, покретачи, ограничења потрошње и броја, контроле извршења, хитне заустављања и квитанте. Не представљајте програм као знак за подршку цене.
\end{itemize}
Тренутни Protocol v1.1.0 Пул таксе, додатне протоколне таксе, и токенски баланси Пул остају регулисани распоређеним уговорима и конфигурацијом, а не овим предложеним вредностима.


\section*{И. Е. Обрађени пример --- предложен мрежни раке}
\addcontentsline{toc}{section}{И. Е. Обрађени пример --- предложен мрежни раке}

Овај пример илуструје предложен модел раке. То није додатни протокол-плата рачунање користи Protocol v1.1.0.

Претпоставимо:

\begin{itemize}[leftmargin=*]
\item У учествујућој групи узима таксу 2.00\%;
\item предложени мрежни грејк добија 20\% ове таксе по групи; и
\item 25\% добитог ракета је прихватљиво и конвертируемо за пријављену употребу у готовности након трошкова.
\end{itemize}
Ефикасна предложена стопа мрежног ракета на маршрутизовану вредност је:

\texttt{rake\_rate = 2.00\% $\times$ 20\% = 0.40\% = 40 bps}

Укупни средства које се могу користити под претпостављеном фактором допустимости од 25\% су:

\texttt{cash\_usable\_rate = 40 bps $\times$ 25\% = 10 bps}

Ако предложен буџет за мрежу има месечни захтев у новчаној номинацији од 150.000 долара и нема другог прихода, илустративна маршрутизована вредност која се захтева по стопи која се може користити у новчаној вредности од 10 бап је:

\texttt{required\_routed\_value = \$150,000 / 0.001 = \$150,000,000 per month}

У овом израчунању се не укључују бруто пошлине покупљања које задржавају покупљања. Додавање таксама у мрежу би двоструко бројело извор.

Реална анализа буџета такође би морала да објави:

\begin{itemize}[leftmargin=*]
\item стварне квитанте за раке и таксе за услугу;
\item прихватљивост актива и трошкове конверзије;
\item учешће у групи и опсег маршрута;
\item усвојили резервне и оперативне циљеве;
\item губици, спорови, исправке и недоступна средства; и
\item Сценарија о осетљивости уместо обећаног раста или повратака.
\end{itemize}
Сваки предложен буџет за приступ таксима или програм ликвидности би остао надолу од усвојена приоритета и могао би бити нулан.


\section*{Ф. Листа за проверу података}
\addcontentsline{toc}{section}{Ф. Листа за проверу података}

\begin{enumerate}[leftmargin=*]
\item \textbf{Историјски докази Sarafu Network}
\begin{itemize}[leftmargin=*]
\item Архивни датум обема суапа, дефиниција активног корисника и средства, инциденти, презентације о откупу, ваучери за старење, потврђено испуњење, ослобођење, мере трајања држања, периоди, искључења и методологија одвојени од тренутних Cosmo-Local Credit метрика.
\item Заштити \href{https://dune.com/grassrootseconomics/sarafu-network}{историјски Даун Аналитика табло}.
\end{itemize}
\item \textbf{Предложен доказ за извршење такса, ако се спроведе}
\begin{itemize}[leftmargin=*]
\item Показати да су било који корак на таксу и реестрски врата интегрисани у извршен пут трансакције, а не да се врши само интерфејсом.
\item Доставите тестове који показују да одговарајућа служба за извршење одбацује пуцање чији примање не доказује прикупљање таксе у складу са усвојеном политиком.
\item Публикујте тестове векторе, случајеве неуспеха, допустима средства и компактне пулове повезане са идентификованом верзијом регистра.
\item Поврзавање одговарајућих \href{https://software.grassecon.org}{опис софтвера}.
\end{itemize}
\item \textbf{Гарантор и пакет покривености}
\begin{itemize}[leftmargin=*]
\item Публикујте допустимост, одговорне стране, финансирање, границе, премије, где је то применимо, изазове, доказе, искључења, овлашћења за доношење одлука, као и примери спора и апелација.
\item Укључите примерни ваучер и обавештења о пулу који разликују одговорност емитента од било које изрично понуђене заштите пула или гаранције треће стране.
\item Заштити \href{https://sarafu.network/reports}{историјски Sarafu Network извештаји архив} и \href{https://youtu.be/5yGaP31bvs0?si=fZ-AMTEXsmVR8Ulv}{презентација историјске године у прегледању} као доказ о породици, а не доказ тренутног покривености CLC или усвајања.
\end{itemize}
\item \textbf{Правни и правни пакет}
\begin{itemize}[leftmargin=*]
\item Документална стратегија надлежности, класе ваучера, политике еквиваленте новчаним средствама, KYC или опције потврде, геофенцирање, откривање за потрошаче, контроле погрешне продаје и разлику између обичног обмена пула и експретног кредитног производа.
\item Везање ефективних \href{https://docs.cosmolocal.credit/sr/governance/terms}{Cosmo-Local Credit Услови коришћења} и чувају замењене верзије кроз контролисане записи.
\end{itemize}
\item \textbf{Усавршавање управљања}
\begin{itemize}[leftmargin=*]
\item Достави процедуре за конфликт интереса и отказ, записи о одлукама и апелацијама, санкције, условни процес за уклањање регистра, границе хитне моћи и примери промена размене и границе у времену.
\end{itemize}
\item \textbf{Извор и пакет за осигурање}
\begin{itemize}[leftmargin=*]
\item Публикујте инвентар извора и лиценце, распоређене адресе и верзије, доступне извештаје о ревизији, инваријанте, произлаз конструкције, овлашћења администратора, контакте инцидента и контролне табле.
\item Поврзавање одговарајућих \href{https://github.com/grassrootseconomics}{Репозиторији за економију}.
\end{itemize}
\end{enumerate}

\end{document}
